\chapter{Problems}


\begin{embeddedproblem}{}
  Use equations~\ref{eq:Exp2} and~\ref{eq:Exp3} to find an IVP for
  which each of the following is the solution:
  \begin{enumerate}[label={\bf{}(\alph*)}]
    \begin{multicols}{3}
    \item $y=2^x$
    \item $y=3^x$
    \item $y=7^x$
    \end{multicols}
  \item $y=a^x$ if $a>0$. (Why do we need to  assume that $a>0$?)
  \end{enumerate}
\end{embeddedproblem}

  \begin{ProblemSection}


  \begin{embeddedproblem}{}
    Show that each of the following differentials is as given:
    \begin{enumerate}[label={\bf{}(\alph*)}]
      \begin{multicols}{2}
      \item $\dx\left(\inverse\sin(2x)\right)=\frac{2}{\sqrt{1-4x^2}}\dx{x}$
      \item $\dx\left(\inverse\sin(2x^2)\right)=\frac{4x}{\sqrt{1-4x^4}}\dx{x}$
      \item $\dx\left(\inverse\sin(2x^3)\right)=\frac{6x^2}{\sqrt{1-4x^6}}\dx{x}$
      \item $\dx\left(\inverse\sin(2x^)\right)=\frac{6x^2}{\sqrt{1-4x^6}}\dx{x}$
      \end{multicols}
      \hrule{}
    \item Now find $d{y}$ when $y=\inverse\sin(2x^n),$ where $n \in\RR.$
    \end{enumerate}  
  \end{embeddedproblem}

  \begin{embeddedproblem}{}
  The sound power level $L$ in decibels for a given sound is given by 
  $$
  L=10\log_{10}\left(\frac{P}{P_0}\right)
  $$
  where $P$ is the sound power of the source measured in watts and
  $P_0$ is the sound reference level taken to be $1$ picowatt or
  $10^{-12}$ watts.
  
  Suppose the sound power of a speaker is $.0001$ watts.  How many
  decibels does this correspond to?

  Suppose at that instant that the sound power of the speaker is being
  raised at a rate of $.001$ watts per second.  How fast is the sound
  power level rising?
\end{embeddedproblem}

  \begin{embeddedproblem}{}
    Show that each of the following differentials is as given:
    \begin{enumerate}[label={\bf{}(\alph*)}]
      \begin{multicols}{2}
      \item $\dx\left(\inverse\sin\left(\frac{x}{\sqrt{x^2+1}}\right)\right)=\frac{1}{\sqrt{x^2+1}}\dx{x}$      
      \item $\dx\left(\inverse\sin\left(\frac{x}{\sqrt{x^2+2}}\right)\right)=\frac{2}{\sqrt{x^2+4}}\dx{x}$      
      \item $\dx\left(\inverse\sin\left(\frac{x}{\sqrt{x^2+5}}\right)\right)=\frac{5}{\sqrt{x^2+25}}\dx{x}$      
      \item $\dx\left(\inverse\sin\left(\frac{x}{\sqrt{x^2+\sqrt{7}}}\right)\right)=\frac{\sqrt{7}}{\sqrt{x^2+7}}\dx{x}$      
      \end{multicols}
      \hrule{}
    \item Now find $d{y}$ when
      $y=\inverse\sin\left(\frac{x}{\sqrt{x^2+a^2}}\right),$ where $n \in\RR.$
    \end{enumerate}
  \end{embeddedproblem}

  \begin{embeddedproblem}{}
    Show that $\dx\left(\inverse\sin(\cos(x))\right)=\dx\left(\inverse\cos(\sin(x))\right)=-1.$
  \end{embeddedproblem}

\end{ProblemSection}

  \begin{embeddedproblem}{}
    Consider a camera $C$ tracking a rocket $R$ traveling vertically
    with a constant speed of $11$  km/second (see below).\\
    \centerline{\includegraphics*[height=3in,width=2in]{../Figures/Prob-TrackingRocket}}
    \begin{enumerate}[label={\bf{(\alph*)}}]
    \item Is the angle $\theta$ increasing at a constant rate?
      Explain.      \item $\cos(y)=x^2+x+3,$ $\left(0, 0\right)$
    \item $\sin(y)=x^2+x+3,$ $\left(0, \frac{\pi}{2}\right)$
    \end{enumerate}
\end{embeddedproblem}

  \begin{embeddedproblem}{}
    Consider a camera $C$ tracking a rocket $R$ traveling vertically
    with a constant speed of $11$  km/second (see below).\\
    \centerline{\includegraphics*[height=3in,width=2in]{../Figures/Prob-TrackingRocket}}
    \begin{enumerate}[label={\bf{(\alph*)}}]
    \item Is the angle $\theta$ increasing at a constant rate?
      Explain.      \item $\cos(y)=x^2+x+3,$ $\left(0, 0\right)$
    \item $\sin(y)=x^2+x+3,$ $\left(0, \frac{\pi}{2}\right)$
    \end{enumerate}
  \end{embeddedproblem}

  \begin{embeddedproblem}{}
    Show that the tangent line to the curve $y=\tan(x)$  at
    $(x_0,y_0)$ is parallel to the tangent line to the curve at
    $(-x_0,-y_0).$
  \end{embeddedproblem}

 

\begin{ProblemSection}
  \begin{embeddedproblem}{}
    Compute each of the following:
    \begin{enumerate}[label={\bf{}(\alph*)}]{}
    \item $\dx(e^{x^2+\sin(x)})$
    \item $\dx(e^{x}\cos(3x)$
    \item $\dx
      \left(
        \frac{3\sqrt{x}+1}{e^x}
      \right)$
    \item $\dx(\tan(e^x))$
    \end{enumerate}{}
  \end{embeddedproblem}

  \begin{embeddedproblem}{}
    For each of the    following, find an equation relating $\dx{x},$
    $\dx{y},$  $\dx{z}.$ Use this to compute $\dfdx{y}{x}$  and
    $\dfdx{y}{t}.$
    \begin{enumerate}[label={\bf{}(\alph*)}]{}
    \item $e^y=x^2+3x-2$
    \item $\sin(xy)=(e^x)^y$
    \item $\tan(x+z)=e^ye^x$
    \item $e^{x^2+y^2+z^2}=2$
    \end{enumerate}
  \end{embeddedproblem}
\end{ProblemSection}

\begin{ProblemSection}
  \begin{embeddedproblem}{}
    Solc the following equations for $x.$
    \begin{enumerate}[label={\bf{}(\alph*)}]{}
    \item $e^x=\frac{1}{\sqrt{e}}$
    \item $e^x=5$
    \item $7e^{x+3}=2$
    \item $\ln(2x-4)=0$
    \item $\ln(2x)-\ln(x+1)=4$
    \item $2\ln(x)=\ln(3x-2)$
    \end{enumerate}
  \end{embeddedproblem}
  \begin{embeddedproblem}{}
    For each of the following find $\dx{y},$ $\dfdx{y}{x},$ and
    $\dfdx{y}{t}.$
    \begin{enumerate}[label={\bf{}(\alph*)}]{}
    \item $y=\ln(2x+2)$
    \item $y=\ln
      \left(
        \frac{(2x+1)^2(3x-2)}{x^2+1}
      \right)$ \hint{Think about properties of logarithms before you
        do this.} 
    \item $\ln(y)=\sin(e^x)$
    \item $\ln(x+y)=\inverse\tan(z)$
    \end{enumerate}
  \end{embeddedproblem}
\end{ProblemSection}

\begin{ProblemSection}
  \begin{embeddedproblem}{}
    Suppose that a point $P$ is moving along the curve
    $y=x^3-x^2+x-1.$ Suppose the horizontal velocity of the point is
    constant.  Would the vertical velocity of the point be constant?
    Justify your answer.
  \end{embeddedproblem}

  \begin{embeddedproblem}{}
    Suppose that a point $P$ starts at the origin and moves along
    the positive $x$-axis with a speed of $1$ unit/sec.  At the same
    time a point $Q$ starts at the point $(1,0)$ and moves up the
    line $x=10$ with a speed of $1$ unit/sec (see below).\\
    \centerline{\includegraphics*[height=2in,width=3.5in]{../Figures/PracProb1}}
    Show that the distance between the two points, $s,$ is decreasing
    $\left(\dfdx{s}{t}<0\right)$ when $y<10-x$ and is increasing
    $\left(\dfdx{s}{t}\right)>0)$ when $y>10-x.$ Use the diagram
    above to describe what this means physically.
  \end{embeddedproblem}
\end{ProblemSection}

\begin{ProblemSection}
  To transform problems into mathematics problems to apply Calculus
  techniques, you need to practice creating equations from word
  problems.  The following exercises provide you the opportunity to
  practice this.

  \begin{embeddedproblem}{}
    Find the slope of the line joining each of the following points to
    an arbitrary point on the curve $ y=x^2+4.$
    \begin{enumerate}[label={\bf{}(\alph*)}]
      \begin{multicols}{3}
      \item $(0,0)$
      \item $(2,3)$
      \item $(2,-3)$
      \item $(-2,3)$
      \item $(-2,-3)$
      \item $(a,b)$
      \end{multicols}
    \end{enumerate}
  \end{embeddedproblem}


  \begin{embeddedproblem}{}
    Two points start at the origin.  One travels along the positive
    $x$-axis with a speed of $3$ units per second.  The other travels
    along the ray $y=x,\ x\ge0$ with a speed of $7$ units per second.  What is
    the distance between these two points at time $t$ seconds.
  \end{embeddedproblem}

  \begin{embeddedproblem}{}
    A marketing company has determined that if a concert ticket costs
    $\$50$ then $20,000$ tickets will sell.  They also determined that
    for each $\$1$ the price is increased, it will cause sales to go
    down by $500$ people.  Let $p$ be the price per ticket.  Find the
    revenue $R$ generated by the concert as a function of $p.$  What
    would the domain of this function be? \hint{Revenue is the
    difference between how much money comes in and how much goes out.}
  \end{embeddedproblem}
   
  \begin{embeddedproblem}{}
    Let $F$ represent temperature in Fahrenheit and $C$ represent
    temperature in Celsius. Use the freezing and boiling points of
    water to derive the formulas to convert Fahrenheit to Celsius and
    from Celcius to Fahrenheit.

    \noindent{}({\sc{}Comment:} If you don't already know the freezing and boiling
    points of water in Fahrenheit and Celsius the first step to
    solving this problem will be to look them up.)
  \end{embeddedproblem}

  \begin{embeddedproblem}{}
    The illumination of an object by a light source is inversely
    proportion to the square of the distance between the source and
    the object.  Consider the following diagram of two street lamps of
    equal strength.\\
  \centerline{\includegraphics*[height=1.2in,width=5in]{../Figures/IllumProb}}

    Write a formula for the illumination $I$ of $P$ by the two street
    lamps.

    \noindent{}({\sc{}Comment:} If you don't know what ``inversely
    proportional'' means the first step to solving this problem will
    be to find out.)
  \end{embeddedproblem}

  \begin{embeddedproblem}{}
    Express the area of an equilateral triangle in terms of its
    perimeter.
  \end{embeddedproblem}
\end{ProblemSection}

\begin{ProblemSection}

   \begin{embeddedproblem}{}
     Given $y={x^2+9x-4}^2,$ compute $\dx{y}$ in three different ways:
     \begin{enumerate}[label={\bf{}(\alph*)}]
       \begin{multicols}{2}
       \item by multiplying before differentiation
       \item by using the Product Rule.
       \item by using the substitution, $z=x^2+9x-4.$
       \end{multicols}
     \item Verify that you get the same answer for each method.
     \end{enumerate}
   \end{embeddedproblem}

   \begin{embeddedproblem}{}\label{prob:1}
     Assume that $x_1, x_2,$ and $x_3$ all depend on $t.$ Use the
     Product Rule to compute $\dfdx{\alpha}{t}$ if:
     \begin{description}
     \item[(a)] $\alpha=x_1x_2.$
     \item[(b)] $\alpha=x_1x_2x_3.$
     \item[(c)] $\alpha=x_1x_2x_3x_4.$
     \item[(d)] $\alpha=x_1x_2x_3x_4x_5.$
     \item[(e)] Find a formula for the derivative of the product of an
       arbitrary number of quantities: $\alpha=x_1x_2x_3\ldots.$
       \hint{Look for a pattern in parts (a) through (d).}
     \end{description}
   \end{embeddedproblem}

    \begin{embeddedproblem}{}
      Let $\alpha=\frac{xy}{z}.$ Use the Product Rule and the Quotient
      Rule in concert to find each of the following:
      \begin{multicols}{2}
        \begin{description}
        \item[(a)] $\dfdx{\alpha}{x}$
        \item[(b)] $\dfdx{\alpha}{y}$
        \item[(c)] $\dfdx{\alpha}{z}$
        \item[(d)] Assume that $x, y,$ and $z$ depend on $t$ and
          compute $\dfdx{\alpha}{t}.$
        \end{description}
      \end{multicols}
    \end{embeddedproblem}

    \begin{embeddedproblem}{}
      Suppose $x=(1+t)^3,$ $y=(1-t)^3.$
      \begin{description}
      \item[(a)] If $z=xy$ show that
        \[
          \dfdx{z}{t} = -6t(1-t^2)^2.
        \]
      \item[(b)] If $z=x/y$ show that
        \[
          \dfdx{z}{t} = \frac{6(1+t)^2}{(1-t)^4}.
        \]
      \end{description}
    \end{embeddedproblem}

    \begin{embeddedproblem}{}
       Find $\dfdx{y}{x}$  and $\dfdx{x}{y}$  for each of the
       following:
       \begin{enumerate}[label={\bf{}(\alph*)}]{}
       \item $\sqrt{x^2+y^2}+xy=5$
       \item $
         \left(
           y+(y+x)^2
         \right)^4=x-y$
       \end{enumerate}
     \end{embeddedproblem}

     \begin{embeddedproblem}{}
       For each of the following curves, find the equation of the
       tangent line at the given point(s).
       \begin{enumerate}[label={\bf{}(\alph*)}]{}
       \item Cardioid: $\left(x^2+y^2\right)^2+4x(x^2+y^2)=4y^2$ at
         $(0,2)$.\\
         \centerline{\includegraphics*[height=2in,width=1.7in]{../Figures/Cardioid}}
       \item Conchoid of Nicomedes: $(y-2)^2(x^2+y^2)=2y^2$ at $(1,1)$
         and $(3,3).$\\
         \centerline{\includegraphics*[height=1.2in,width=2in]{../Figures/ConchoidOfNicomedes}}
         Note: Nicomedes (circa 280 BC - 210 BC) used this curve to
         trisect an angle, a construction which is impossible to
         accomplish with a straight edge and compass alone.  
       \item Cassini oval: $(x^2+y^2+1)^2-4xx^2=\frac32$ at
         $\left(\frac{\sqrt[4]{2}}{2},\frac{\sqrt[4]{2}}{2}\right)$\\
         \centerline{\includegraphics*[height=1.2in,width=2in]{../Figures/CassiniOval}}
         Note: Giovanni Cassini (1625-1712) studied Cassini Ovals in
         astronomy.  Whereas ellipses are points the sum of whose
         distances from two fixed foci is constant, Cassini ovals are
         where the product of the distances is constant.  In honor of
         his accomplishments in astronomy, the Cassini space probe was
         launched in 1997 and became the first probe to orbit Saturn. 
       \end{enumerate}
     \end{embeddedproblem}
     \begin{embeddedproblem}{}
       Suppose $P(t)= \left(x \left (t \right ) ,y \left (t \right )
       \right)$  is a point moving in the plane along the following
       curves.
       \begin{enumerate}[label={\bf{}(\alph*)}]{}
       \item Find an equation that relates its vertical speed,
         $\dfdx{y}{t}$  to it horizontal speed, $\dfdx{x}{t}$.
         \begin{enumerate}[label={\bf{}(\roman*)}]{}
         \item $\sqrt {{x} ^ {2} + {y} ^ {2}} +xy= {\left (x+y \right )} ^    {3}$
         \item ${\left (x-3 \right )} ^ {2} + {\left (y+5 \right )} ^ {2} =  {r} ^ {2}$
         \item ${\left ({x} ^ {2} -3y \right )} ^ {3} =2x+y$
         \end{enumerate}
       \item For each of the curves in Part a, suppose
         $\dfdx{x}{t}=1\frac{\text{unit}}{\text{second}}$. Find the
         speed, $\dfdx{s}{t}=\sqrt{\left(\dfdx{x}{t}\right)^2+
           \left(\dfdx{y}{t}\right)^2}$   in terms of $x$ and $y$.
       \end{enumerate}

     \end{embeddedproblem}
   % \begin{embeddedproblem}{}
    %   Use the Product Rule to compute $\dx y:$
    %   \begin{multicols}{2}
    %     \begin{description}
    %     \item[(a)] $y=x\cdot x$
    %     \item[(b)] $y=x^4(x-1)^3$
    %     \end{description}
    %   \end{multicols}
    % \end{embeddedproblem}

  \begin{embeddedproblem}{}
    For each of the following find all solutions of the equation
    $\dfdx{y}{x} =0.$
    \begin{multicols}{3}
      \begin{enumerate}[label={\bf{}(\alph*)}]
      \item $ y=x^3(x+2)^4$
      \item $y=x^2(x+2)^3$
      \item $y=\inverse{x}(x+2)^{-7}$
      \item $y=x^a(x+2)^b$
      \item $y=x^a(x+c)^b$
      \end{enumerate}
    \end{multicols}
  \end{embeddedproblem}
  \begin{embeddedproblem}{}
    For each of the following find all solutions of the equation
    $\dfdx{y}{x} =0.$
    \begin{multicols}{3}
      \begin{enumerate}[label={\bf{}(\alph*)}]
      \item $ y=\frac{x^3}{(x+2)^4}$
      \item $ y=\frac{x }{(x+2)}$
      \item $ y=\frac{x^2}{(x+2)^2}$
      \item $ y=\frac{x }{(x+2)^2}$
      \item $ y=\frac{x^2}{(x+2) }$
      \item $ y=\frac{x^3}{(x+2)^5}$
      \item $ y=\frac{x^a}{(x+2)^b}$
      \end{enumerate}
    \end{multicols}
  \end{embeddedproblem}
  \begin{embeddedproblem}{}
    For each of the following find all solutions of the equation
    $\dfdx{y}{x} =0.$
    \begin{multicols}{2}
      \begin{enumerate}[label={\bf{}(\alph*)}]
      \item $ y=\sqrt{3x}\inverse{(x-1)}$
      \item $ y=\frac{\sqrt{3x}}{x-1}$
      \item $ y=(3x-1)^7$
      \item $ y=3(x-1)^7$
      \item $ y=(5x-4)^2(2x+1)^3$
      \item $ y=\sqrt{2x-1}+\sqrt[3]{2x-1}$
      \item $ y=\left(2\sqrt{x}-1\right)+\left(2\sqrt[3]{x}-1\right)$
      \item $ y=\sqrt{1-x^2}$
      \item $ y=x\sqrt{1-x^2}$
      \item $ y=\frac{1}{\sqrt{1-x^2}}$
      \item $ y=\frac{x}{\sqrt{1-x^2}}$
      \item $ y=\frac{1}{\sqrt{1-x^3}}$
      \item $ y=\frac{x}{\sqrt{1-x^3}}$
      \item $ y=\frac{1}{\sqrt{1-x^4}}$
      \item $ y=\frac{x}{\sqrt{1-x^4}}$
      \end{enumerate}
    \end{multicols}
  \end{embeddedproblem}

  \begin{embeddedproblem}{}
    Find $\dfdx{y}{x}$ and $\dfdx{x}{y}$ for each of the following.
    \begin{multicols}{2}
      \begin{description}
      \item[(a)] $x^2-4xy+y^2=4$
      \item[(b)] $2x^2+xy-y^2=3x$
      \item[(c)] $2x^2+xy=y^2+3x$
      \item[(d)] $y\cos(x) = x^2+y^2$
      \item[(e)] $x\sec(y) = x^2-y^2$
      \item[(f)] $(x^2+y^2)^2 = x^2-y^2$
      \end{description}
    \end{multicols}
  \end{embeddedproblem}

  \begin{embeddedproblem}{}
    For each of the following show that $\dfdx{y}{x}$ is as given.
    \begin{enumerate}[label={\bf{}(\alph*)}]
      \begin{multicols}{2}
      \item $\displaystyle y=\frac{x^2-1}{x-1} = 1$
      \item $\displaystyle y=\frac{x^3-1}{x-1} = 2x+1$
      \item $\displaystyle y=\frac{x^4-1}{x-1} = 3x^2+2x+1$
      \item $\displaystyle y=\frac{x^5-1}{x-1} = 4x^3+3x^2+2x+1$
      \item $\displaystyle y=\frac{x^6-1}{x-1} = 5x^4+ 4x^3+3x^2+2x+1$
      \end{multicols}
    \item Find $\displaystyle y=\frac{x^n-1}{x-1},$ where $n$ is a
      positive integer.
    \item Confirm that the formula you found in part (f) works for
      parts (a), (b), (c), (d), and (e).
    \end{enumerate}
  \end{embeddedproblem}

  \begin{embeddedproblem}{}
    For this problem $x_1, x_2, x_3, \ldots, x_n$ are all
    constants. The variable is $x,$ as usual.
    \begin{enumerate}[label={\bf{}(\alph*)}]
    \item Show that:
      $\dx (x-x_1)(x-x_2) = \left[(x-x_1)+(x-x_2)\right]\dx x$
    \item Show that:
      \[\dx (x-x_1)(x-x_2)(x-x_3) = \left[(x-x_2)(x-x_3) +
          (x-x_1)(x-x_3) + (x-x_1)(x-x_2)\right]\dx x\]
    \item Show that:
      \begin{multline*}
        \dx (x-x_1)(x-x_2)(x-x_3)(x-x_4)\\ = [(x-x_2)(x-x_3)(x-x_4) +
        (x-x_1)(x-x_3)(x-x_4)\\ + (x-x_1)(x-x_2)(x-x_4) +
        (x-x_1)(x-x_2)(x-x_3)]\dx x
      \end{multline*}
    \item Guess what the differential of
      \( (x-x_1)(x-x_2)(x-x_3)(x-x_4)(x-x_5) \) is and show that your
      guess is correct.
    \item Guess what the differential of
      \( (x-x_1)(x-x_2)(x-x_3)\ldots(x-x_n) \) is. Can you show that
      your guess is correct?
    \end{enumerate}
  \end{embeddedproblem}

  \begin{embeddedproblem}{}
    Boyle's Law says that if a gas is held at a constant temperature
    the the product of its pressure ($P$) and it volume ($V$) is
    constant.
    \begin{enumerate}[label={\bf{}(\alph*)}]
    \item Show that if the temperature is constant then
      \[\dfdx{P}{V} = -\frac{P}{V}.\]
    \item Suppose the volume of a gas is increasing at $2$ cubic
      inches per second. Find the rate of change of the pressure with
      respect to time, $\dfdx{P}{t}.$ What is $\dfdx{P}{t}$ if the
      volume is increasing at $3$ cubic inches per second? $4$ cubic
      inches per second? $5?$ What if it is increasing at $n$ cubic
      inches per second?
    \end{enumerate}
  \end{embeddedproblem}

  \begin{embeddedproblem}{}
    Suppose the price, $P,$ of a commodity is increasing at a rate of
    5\% and the quantity sold, $Q,$ is decreasing at rate of 7\%.  Is
    the revenue, $R,$ generated by the sale of this commodity
    increasing or decreasing?  What is the percent rate of change?
  \end{embeddedproblem}

  \begin{embeddedproblem}{}
    The curve given by the equation
    $$
    (x^2+x^2)^2=x^2-y^2
    $$
    is an example of a curve called a lemniscate.  The graph is below.\\
    \centerline{\includegraphics*[height=2in,width=4.5in]{../Figures/LemniscateProb}}
    Find the points where the tangent lines are horizontal.
  \end{embeddedproblem}

  \begin{embeddedproblem}{}
    \begin{enumerate}[label={\bf{}(\alph*)}]
    \item A point moves in the plane so that its coordinates at time
      $t$ are given by $x=t$ $y=t^2.$
      \begin{enumerate}[label={\bf{}(\roman*)}]
        % \begin{multicols}{2}
      \item Find the slope of the tangent line at times $t=0, 1, 2.$
      \item Find the speed of the point at times $t=0,1,2.$
        % \end{multicols}
      \end{enumerate}
    \item A point moves in the plane so that its coordinates at time
      $t$ are given by $x=t^2$ $y=t^4.$
      \begin{enumerate}[label={\bf{}(\roman*)}]
        % \begin{multicols}{2}
      \item Find the slope of the tangent line at times $t=0, 1, 2.$
      \item Find the speed of the point at times $t=0,1,2.$
        % \end{multicols}
      \end{enumerate}
    \item A point moves in the plane so that its coordinates at time
      $t$ are given by $x=t^3+t+1,$ $y=t^4+7t-3.$
      \begin{enumerate}[label={\bf{}(\roman*)}]
        % \begin{multicols}{2}
      \item Find the slope of the tangent line at times $t=0, 1, 2.$
      \item Find the speed of the point at times $t=0,1,2.$
      \end{enumerate}
    \end{enumerate}
  \end{embeddedproblem}


  \begin{embeddedproblem}
    Suppose a point is moving along the curve $x^2+y^3=\sqrt{xy}$.
    Find an equation relating the horizontal and vertical
    accelerations.
  \end{embeddedproblem}
  

  \begin{embeddedproblem}{}
    Suppose $y=y(x)$ and $x=x(t).$ Show that
    $$
    \dfdxn{y}{t}{2}=\dfdx{y}{x} \left( \dfdxn{x}{t}{2} \right)
    +\dfdxn{y}{x}{2} \left( \dfdx{x}{t} \right)^2
    $$
  \end{embeddedproblem}

  \begin{embeddedproblem}{}
    \begin{multicols}{2}[Compute $\dfdxn{y}{x}{2}$ at the given point]
      \begin{enumerate}
      \item $x^2+y^2=2$ at $(1,1)$
      \item $x^2-y^2=16$ at $(5,3)$
      \item $x^3y+xy^3=2$ at $(1,1)$
      \end{enumerate}
    \end{multicols}
  \end{embeddedproblem}

  \begin{embeddedproblem}{}
      \begin{enumerate}[label={\bf{}(\alph*)}]
      \item Compute $\dx{y}$ for each of the following:
        \begin{multicols}{2}
          \begin{enumerate}[label={\bf{}(\roman*)}]
          \item $ y=\frac{x^2-2x+4}{x}$
          \item $ y=x^2-\sqrt{x}$
          \item $ y=\frac{x}{x+1}$
          \item
            $ y=\left(\sqrt[3]{x}\right)^2-\left(\sqrt{x}\right)^3$
          \item
            $ y=\left(\sqrt[3]{x}\right)^3-\left(\sqrt{x}\right)^2$
          \item
            $
            y=\left(\frac{1}{x^2}+\frac{1}{x^3}\right)\left(x^2+x^3\right)$
          \item $ y=\left(x^2+x^3\right)\left(x^{-2}+x^{-3}\right)$
          \item
            $
            y=\inverse{\left(x^2+x^3\right)}\left(x^{-2}+x^{-3}\right)$
          \item $ y=\frac{x^2-1}{x-2}$
          \item $ y=\frac{x^3-1}{x-3}$
          \item $ x^2+y^2=1$
          \end{enumerate}
        \end{multicols}
      \item Redo each of the problems in part (a), but this time
        compute $\dfdx{y}{x}$, the ratio of the differentials $\dx{y}$ and
        $\dx{x}.$
      \end{enumerate}
    \end{embeddedproblem}

    \begin{embeddedproblem}{}
      Compute each of  the following:
      \begin{enumerate}[label={\bf{}(\alph*)}]
        \begin{multicols}{2}
        \item $\dx 
          \left(
            \frac{x^3+2x^2}{\inverse{x}+7x^{15}}
          \right)$
        \item $\dx
          \left(
            \frac{(3x^2+2x-1)(4x^8-x^{-9}+2)}{\sqrt{x}+\frac{2}{\sqrt{x}}}
          \right)$
        \end{multicols}
      \end{enumerate}
    \end{embeddedproblem}


      %       What we've done so far is just the beginning for the rules for
      %       Calculus and equivalent formulations of what we've done were known
      %       before Newton and Leibniz.  Recall the title of Leibniz' first paper
      %       on differential Calculus:
      %       \begin{center}
      %       \begin{minipage}{0.85\linewidth}
      %       \emph{A New Method for Maxima and Minima, as Well as Tangents, Which is
      %       Impeded Neither by Fractional nor Irrational Quantities, and a
      %       Remarkable Type of Calculus for This.}
      %       \end{minipage}
      %       \end{center}
      %       A reasonable question to ask
      %       is ``What do these have to do with tangents, maxima, or minima?''
      %       Another is ``What does he mean by is Impeded neither Fractional or
      %       irrational quantities?''  These are good questions and we'll address
      %       both of these in the next section as we expand our rules to
      %       incorporate the rest of the differentiation rules and see where slopes
      %       come into play.  This generality is why Newton and Leibniz get credit
      %       for inventing Calculus instead if their predecessors.

      %       Recall that the title of Leibniz's first paper on Calculus made
      %       reference to ``Maxima and Minima, as Well as Tangents.'' A reasonable
      %       question to  ask is ``What do these rules have to do with tangents,
      %       maxima or minima?''  Good question!  Our gateway to this
      %       comes in the form of an old friend, the slope of a line. 



  
  

  % \begin{embeddedproblem}{}
  %   Suppose $p(x)$ is a quadratic polynomial with a double root $r.$
  %   Thus $p(x)=a(x-r)^2$
  %   \begin{enumerate}
  %   \item Expand $p(x)$ and use this to form the Hudde polynomial
  %     $p^*(x).$ Show that $p^*(r)=0$ (so that $r$ is a root of
  %     $p^*(x)).$
  %   \item Show that if $r$ is a double root of the polynomial $p(x)$
  %     then it is a root of $\dfdx{p}{x}.$ [Hint: If $r$ is a double
  %     root of $p(x),$ then $p(x)=(x-r)^2 q(x)$ for some polynomial
  %     $q(x).$]
  %   \item Show that the Hudde polynomial
  %     $p^*(x)=ap(x)+bx\dfdx{p}{x}$ and use this to prove Hudde's
  %     Rule.
  %   \end{enumerate}
  % \end{embeddedproblem}



      %       \begin{enumerate}
      %       \item $\dfdx{ (\sin x)}{x} = \cos x $
      %       \item $\dfdx{ (\cos x)}{x} = -\sin x $
      %       \item $\dfdx{ (\tan x)}{x} = \sec^2 x $
      %       \item $\dfdx{ (\cot x)}{x} = -\csc^2 x $
      %       \item $\dfdx{ (\sec x)}{x} = \sec x\tan x $
      %       \item $\dfdx{ (\csc x)}{x} = -\csc x\cot x $
      %       \end{enumerate}

      %       \begin{enumerate}
      %       \item $\displaystyle\dfdx{ (\inverse\sin x)}{x} =  \frac{1}{\sqrt{1-x^2}} $
      %       \item $\displaystyle\dfdx{ (\inverse\cos x)}{x} = \frac{-1}{\sqrt{1-x^2}}  $
      %       \item $\displaystyle\dfdx{ (\inverse\tan x)}{x} = \frac{1}{1+ x^2} $
      %       \item $\displaystyle\dfdx{ (\inverse\cot x)}{x} = \frac{-1}{1+ x^2}$
      %       \item $\displaystyle\dfdx{ (\inverse\sec x)}{x} = \frac{1}{\abs{x}\sqrt{x^2-1}} $
      %       \item $\displaystyle\dfdx{ (\inverse\csc x)}{x} = \frac{-1}{\abs{x}\sqrt{x^2-1}} $
      %       \end{enumerate}


% \section{Differentials, Velocity, and the Slope of a Line}
% \label{sec:diff-slope-line}
% \markright{{\sc{}Differentials, Velocity, Slope}}

% Now that we have the idea of differentials there is an obvious
% question that deserves to be addressed: ``What's the point of all of
% this?''  The differentials themselves aren't really so important
% though we will make use of them later. But the ratios of differentials
% turn out to be very important\footnote{\textcolor{red}{Need to add
%     some discusion of $\dfdx{x}{t}$ as a velocity.}}.  For example, recall that in
% Section~\ref{sec:ferm-meth-aedeq} we used Fermat's method show that
% the slope of the tangent line to the curve $y=x^2$ at the point
% $(x, x^2)$ is given by $2x$ and that this was consistent with the
% graph.  Using this \emph{differential} calculus, we can
% obtain these results more directly.

% For example, from our efforts in the previous section we know that if
% $y=x^2,$ then $\d{y}=2x\dx{x}.$ If we divide both sides of this last
% formula by $d{x}$ we recapture Fermat's result: $\dfdx{y}{x}=2x.$

% Is this the slope of the tangent line?  Think about it for a moment.
% If $(x_0,y_0)$ and $(x_1,y_1)$ are two points on the tangent line then
% the slope is the change in $y$ divided by the change in $x:$
% $\frac{y_1-y_0}{x_1-x_0}.$ If we think of these points as being
% infinitely close together, then the slope is given by $\dfdx{y}{x},$
% so the slope of the line tangent to a curve at a point is given by the
% ratio: $\dfdx{y}{x}.$

% We can also think of $\dfdx{y}{x},$ loosely speaking, the slope of the
% curve \emph{itself} at a point. After all, if $\dx{y}$ $\dx{x}$ are
% \emph{infinitely small} then we're really talking about the slope of
% the curve at a single point, aren't we?
% \centerline{\includegraphics*[height=1in,width=2in]{../Figures/parabola1}}
% This characterization should trouble you. If it doesn't, ask yourself
% this question: If we're talking about a \emph{single} point $(x,y)$
% then what are $\dx{x}$ and $\dx{y}?$ Aren't they both just
% zero\footnote{This is a bit of a conundrum, but not one we need to
%   resolve right now. We'll come back to it later.}?

% % let's look at the tangent line to the parabola at the point
% % $(x,x^2).$ Recall from section~\ref{sec:descartes-normals} that Fermat
% % determined that the slope of this tangent line is given by $2x.$
% % Dividing through by $\dx{x}$ gives
% % $$
% % \dfdx{y}{x}=\dfdx{(x^2)}{x}=\frac{2x\dx{x}}{\dx x}=2x.
% % $$
% % The only difference between $\dfdx{y}{x}=2x$ and $\frac{\Delta
% %   y}{\Delta x}=m$ is that $\d{y}$ and $\dx{x}$ are infinitesimals,
% % whereas $\Delta y$ and $\Delta x$ are finite.

% To see that we (the authors) haven't completely lost touch with
% reality, let's get infinitely close to the curve at the generic point
% $(x,x^2)$ to see what is happening. Well, if we think we can actually
% get \underline{infinitely close} then we really have lost touch with
% reality. Let's get really, really close.

% \centerline{\includegraphics*[height=2in,width=5in]{../Figures/parabola-zoom}}

% As you can see from the sequence of images above, as we zoom in, it
% becomes harder to tell the difference between the curve and the
% tangent line.  If we could zoom in infinitely close, then the curve
% and the tangent line would become indistinguishable. Moreover, since
% the the expression $\dfdx{y}{x}$ is the (infinitely small) change in
% $y$ divided by the (infinitely small) change in $x$ it is clearly
% \emph{both} the slope of the curve at the point of tangency \emph{and}
% the slope of the tangent line at the same point. We can safely think
% of it in either mode, and it will be helpful to switch between modes
% as needed.

% In problem~\ref{prob:ferm-meth-x3}, we saw that it is very hard to use
% Fermat's method to determine the slope of the tangent line to the
% curve $y=x^3$ at the point $(x,x^3).$ However the Power Rule, as
% stated by Leibniz, gives the answer easily: $\dx{y}=d(x^3 )=3x^2 \dx x$
% or $\dfdx{y}{x}=3x^2.$
% That this last expression is not a constant reflects the fact that the
% slope of the curve $y=x^3$ changes at every value of $x.$

% Leibniz's claim that his ``New Method for Tangents,'' which works ``as
% by magic,'' seems to be justified.

\begin{embeddedproblem}{}
  You are probably aware of the fact, from elementary geometry, that
  the line tangent to to a circle is always perpendicular to the
  radius at the point of tangency. Recall that Descartes used this
  fact to derive his Method of Normals.

  \begin{enumerate}[label={\bf{}(\alph*)}]
  \item  The equation of the circle centered at the origin with radius
    equal to $5$ is  $$x^2+y^2=25.$$ Find an equation relating
    $\dx{y}$ and $\dx{x}.$
  \item Compute $\dfdx{y}{x}$ and compare this to the slope of the
    radius connecting the origin to each of the following points:
    \begin{enumerate}[label={\bf{}(\roman*)}]    
    \item  $(3,4)$ 
    \item $(2, \sqrt{21}$
    \end{enumerate}
  \item The equation of the circle centered at the point $(1,2)$ with
    a radius of $6$ is $$(x-1)^2+(y-2)^2=36.$$ Find an equation relating
    $\dx{y}$ and $\dx{x}.$
  \item Compute $\dfdx{y}{x}$ and compare this to the slope of the
    radius connecting the point $(1,2)$ to each of the following points:
    \begin{enumerate}[label={\bf{}(\roman*)}]    
    \item $(0,8)$
    \item $(0,-4)$
    \item $(7,0)$
    \item $(5,0)$
    \end{enumerate}
  \item The equation of the circle centered at the point $(h,k)$ with
    a radius of $r$ is $$(x-h)^2+(y-k)^2=r^2.$$ Find an equation relating
    $\dx{y}$ and $\dx{x}.$
  \item Assume that the point $(a,b)$ is on the circle above. Compute $\dfdx{y}{x}$ and compare this to the slope of the
    radius connecting the point $(h,k)$ to the point $(a,b).$
  \end{enumerate}
\end{embeddedproblem}
  \end{ProblemSection}

  
\begin{ProblemSection}

\end{ProblemSection}


\begin{ProblemSection}
\end{ProblemSection}

\begin{ProblemSection}
\begin{ProficiencyDrill}
  Find the equation of the tangent line(s) for the given values of
  $x$.
  \begin{enumerate}[label={\bf{}(\roman*)}]
    \begin{multicols}{2}
      \item
        $y = \abs{x^3}$\\ at the point $ x=0$
      \item
        $y = -\abs{x^3}$\\ at the point $ x=1$
      \item
        $y = x^3-\dfrac1x$\\ at the point $x=1$
      \item
        $y =
        \begin{cases}
          x^2& x< 2\\
          4 & x=2\\
          \frac13(x^3+4)& x>2
        \end{cases}$\\ at the point $x=2$
      \item
        $y=\sin(\sin(x))$\\ at the point $(\pi,0)$
      \item
        $y=\sin^2(\cos(x))$\\ at the point $(\pi/2,0)$
      \item
        $y=x\sqrt[3]{x-8}$\\ at the point $(0,0)$
      \item
        $y=x\sqrt[3]{x-8}$\\ at the point $x=6$
      \item
        $y=x\sqrt[3]{x-8}$\\ at the point $x=8$
      \item
        $y=\sqrt[3]{x}{(x-8)}$\\ at the point $x=8$
      \item
        $y=\sqrt[3]{x}{(x-8)}$\\ at the point $(0,0)$
      \item
        $y= \dfrac{1}{2+\sin(x)}$\\ at the point $(0,1/2)$
      \item
        $y= \dfrac{1}{2+\sin(x)}$\\ at the point $(\pi/2,1/3)$
      \item
        $y= \dfrac{x}{e^x}$\\ at the point $(0,0)$
      \item
        $y= \dfrac{x}{e^x}$\\ at the point $\left(1,1/e\right)$
      \item
        $y= \tan(x)$\\ at the point $(\pi/4,1)$
      \item
        $y= \cot(x)$\\ at the point $(\pi/4,1)$
      \item
        $y= \dfrac{e^x}{4x^2+3x-5}$\\ at the point
        $\left(0,\dfrac{-1}{5}\right)$
      \item
        $y= \dfrac{e^x}{4x^2+3x-5}$\\ at the point $(1,e/2)$
      \item
        $y= \sin(x)(1+\cos(x))$\\ at the point $x=\pi/3$
      \item
        $y= \sin(x)(1+\cos(x))$\\ at the point $x=-\pi/3$
      \item
        $y= \sin(x)\tan(x)+\cos(x)-\sec(x)$\\ at the point $(0,0)$
      \item
        $y= \sin(x)\tan(x)+\cos(x)-\sec(x)$\\ at the point $x=\pi/4$
      \item
        $y= x\sin(x)+\cos(x)$\\ at the point $(0,1)$
      \item
        $y= x\sin(x)+\cos(x)$\\ at the point $(\pi,-1)$
    \end{multicols}
  \end{enumerate}
\end{ProficiencyDrill}
\end{ProblemSection}

\begin{ProblemSection}

    \begin{embeddedproblem}{}
    Recall that we used Fermat's Method to compute the slope of the
    tangent line to $y= {x} ^ {3} +2 {x} ^ {2}$.  Use the rules of
    calculus to compute $\dfdx{y}{x}$  and compare
    this to the answer Fermat obtained without calculus.  Which method
    was easier? 
  \end{embeddedproblem}

  \begin{embeddedproblem}{}
    Suppose $g \left (x \right ) =f \left (x \right ) +b$ where is
    some constant.  Show that $\dfdx{y}{x}=\dfdx{g}{x}$.  What does
    this say about the tangent lines to the curves $y=f(x)$ and
    $y=g(x)$?  Does this make sense geometrically?  Explain.
  \end{embeddedproblem}

    \begin{embeddedproblem}{}
    Suppose a point $P$ is moving along the following curves so that
    its horizontal velocity is
    $\dfdx{x}{t}= 2\frac{\text{units}}{\text{second}}.$ Find its
    vertical velocity at the instant the point is located at the
    indicated position.
    \begin{enumerate}[label={\bf{}(\alph*)}]
    \item $x^2+y^2=2$ at $(1,1)$
    \item $y=y= {x} ^ {3} +5 {x} ^ {2} -9 $ at $(0,9)$
    \item $x\sqrt{y^3}=y^2+\frac{1}{\sqrt[3]{x}}$ at $(x,y)$
    \end{enumerate}
  \end{embeddedproblem}
  \begin{embeddedproblem}{}
    For each curve, compute $\dfdx{y}{x}$ and use this to find the equation of
    the tangent line to the curve at the indicated point.
    \begin{enumerate}[label={\bf{}(\alph*)}]
    \item  $ \frac{x^2}{27}+\frac{y^2}{9}=1,$ $ (3,\sqrt{6})$
    \item  $    xy=1,$ $(-2,-1/2)$
    \item  $    y=\sqrt{x},$ $ (4,2)$
    \item  $    y=2x^3-3x+5,$ $ (-1,6)$
    \item  $    y=\frac{x^3-2x^2+3}{2x-1},$ $ (1,2)$
      \textcolor{red}{This part of this problem is misplaced. Needs to be moved.}
    \end{enumerate}
  \end{embeddedproblem}




  \begin{embeddedproblem}{}
    Find an equation of the line tangent to each of the curves below at
    the points indicated.   To check your calculations graph both the curve and the tangent
    line. You should be able to tell visually whether you are correct.
    \begin{enumerate}[label={\bf{}(\alph*)}]
    \item $y=x^2$
      \begin{enumerate}[label={\bf{}(\roman*)}]
        \begin{multicols}{3}
        \item $(0,0)$
        \item $(1,1)$
        \item $(-1,1)$
        \item $(2,4)$
        \item $(-2,4)$
        \item $(5,25)$
        \end{multicols}
      \item Now find an equation of the tangent line at $(a,a^2)$
        where $a$ is any real number.
      \end{enumerate}
      \hrule{}
    \item $y=x^3$
      \begin{enumerate}[label={\bf{}(\roman*)}]
        \begin{multicols}{3}
        \item $(0,0)$
        \item $(1,1)$
        \item $(-1,-1)$
        \item $(2,8)$
        \item $(-2,-8)$
        \item $(5,125)$
        \end{multicols}
      \item Now find the equation of the tangent line at $(a,a^3)$
        where $a$ is any real number.
      \end{enumerate}
      \hrule{}
    \item $y=x^3+x^2$
      \begin{enumerate}[label={\bf{}(\roman*)}]
        \begin{multicols}{3}
        \item $(0,0)$
        \item $(1,2)$
        \item $(-1,0)$
        \item $(2,12)$
        \item $(-2,-4)$
        \item $(5,150)$
        \end{multicols}
      \item Now find the equation of the tangent line at
        $(a,a^3+a^2)$ where $a$ is any real number.
      \end{enumerate}
      \hrule{}
    \item $y=x^4-3x^2+5$
      \begin{enumerate}[label={\bf{}(\roman*)}]
        \begin{multicols}{3}
        \item $(0,5)$
        \item $(1,3)$
        \item $(-1,3)$
        \item $(2,9)$
        \item $(-2,9)$
        \item $(5,555)$
        \end{multicols}
      \item Now find the equation of the tangent line at
        $(a,a^4-3a^2+5)$ where $a$ is any real number.
      \end{enumerate}
      \hrule{}
    \item $y=2x^3+2x-4$
      \begin{enumerate}[label={\bf{}(\roman*)}]
        \begin{multicols}{3}
        \item $(0,-4)$
        \item $(1,0)$
        \item $(-1,-4)$
        \item $(2,16)$
        \item $(-2,-24)$
        \item $(5,256)$
        \end{multicols}
      \item Now find the equation of the tangent line at
        $(a,2a^3+2a-4)$ where $a$ is any real number.
      \end{enumerate}
      \hrule{}
    \end{enumerate}
  \end{embeddedproblem}


  \begin{embeddedproblem}{}
    Find an equation of the line tangent to each of the curves
    below at the points indicated.   To check your calculations graph both the curve and the tangent
    line. You should be able to tell visually whether you are correct.

    \begin{enumerate}[label={\bf{}(\alph*)}]
    \item $y=\frac{x+1}{x-2}$
      \begin{enumerate}[label={\bf{}(\roman*)}]
        \begin{multicols}{3}
        \item $(0,-1/2)$
        \item $(1,-2)$
        \item $(-1,0)$
        \item $(3,1)$
        \item $(-3,-2/5)$
        \item $(a,\frac{a+1}{a-2})$ where $a$ is any real number.
        \end{multicols}
      \end{enumerate}
    \item $y=\frac{2x+1}{3x-2}$
      \begin{enumerate}[label={\bf{}(\roman*)}]
        \begin{multicols}{3}
        \item $(0,-1/2)$
        \item $(1,3)$
        \item $(-1,1/5)$
        \item $(3,1)$
        \item $(-3,5/7)$
        \item $(a,\frac{2a+1}{3a-2})$ where $a$ is any real number.
        \end{multicols}
      \end{enumerate}
    \item $y=\frac{ax+b}{cx+d}$ where $a, b, c,$ and $d$ are fixed,
      but unspecified real numbers.
      \begin{enumerate}[label={\bf{}(\roman*)}]
        \begin{multicols}{3}
        \item $\left(1,\frac{a+b}{c+d}\right)$
        \item $\left(-1,\frac{-a+b}{-c+d}\right)$
        \item $\left(2,\frac{2a+b}{2c+d}\right)$
        \item $\left(-2,\frac{-2a+b}{-2c+d}\right)$
        \item $\left(\alpha,\frac{\alpha a+b}{\alpha c+d}\right)$
          where $\alpha$ is any real number.
        \end{multicols}
      \end{enumerate}

    \item $y=\frac{x^2-2}{x+1}$ 
      \begin{enumerate}[label={\bf{}(\roman*)}]
        \begin{multicols}{3}
        \item $(0,-2)$
        \item $(1,-1/2)$
        \item $(2,2/3)$
        \item $(-2,-2)$
        \item $\left(a,\frac{x^2-2}{x+1}\right)$ where $a$ is any real
          number. 
        \end{multicols}
      \end{enumerate}
    \end{enumerate}
  \end{embeddedproblem}



  

\end{ProblemSection}

% \begin{ProblemSection}
%   \begin{embeddedproblem}{}
%     Find an equation of the line tangent to each of the curves
%     below at the points indicated. To check your calculations graph
%     both the curve and the tangent line. You should be able to tell
%     visually whether you are correct.
%     \begin{enumerate}[label={\bf{}(\alph*)}]
%     \item $y=x^2$ at
%       \begin{enumerate}[label={\bf{}(\roman*)}]
%         \begin{multicols}{2}
%         \item $(1,1)$
%         \item $(-1,1)$
%         \item $(2,4)$
%         \item $(-2,4)$
%         \item $(a,a^2)$
%         \end{multicols}
%       \end{enumerate}
%     \item $y=\frac{2x+1}{3x-2}$
%       \begin{enumerate}[label={\bf{}(\roman*)}]
%         \begin{multicols}{3}
%         \item $(0,-1/2)$
%         \item $(1,3)$
%         \item $(-1,1/5)$
%         \item $(3,1)$
%         \item $(-3,5/7)$
%         \item $(a,\frac{2a+1}{3a-2})$ where $a$ is any real number.
%         \end{multicols}
%       \end{enumerate}
%     \end{enumerate}
%   \end{embeddedproblem}

% \end{ProblemSection}

  \begin{ProblemSection}
    \begin{embeddedproblem}{}
      Show that the following families of curves are orthogonal where
      they intersect.  Sketch a few curves from each family to see if
      they look perpendicular.
      \begin{enumerate}[label={\bf{}(\alph*)}]{}
      \item $xy=a$ and $ {x}^{2} - {y}^{2} =b$
      \item $y=c {x} ^ {3}$ and ${x} ^ {2} +3 {y} ^ {2}
        =a$
      \item $y= c {x} ^ {n}$ and $ {x} ^ {2} +n {y} ^ {2}
        =a$ where $n$ is a fixed rational number (fraction).
      \item ${x} ^ {2} + {y} ^ {2} =ax$ and $          {x} ^ {2} + {y} ^ {2}   =by$
      \end{enumerate}
    \end{embeddedproblem}
  
  \end{ProblemSection}

\begin{ProblemSection}

\begin{embeddedproblem}{}
    Suppose\footnote{\textcolor{red}{This entire problem section
        should probably be relocated. ecb; 3/8/20}} that \(3x-4y+7z=2.\)
        \begin{description}
        \item[(a)] Assuming that $y$ and $z$ depend on $x,$ compute
          $\dfdx{y}{x}$ and $\dfdx{z}{x}.$
        \item[(b)] Assuming that $x$ and $y$ depend on $z,$ compute
          $\dfdx{x}{z}$ and $\dfdx{y}{z}.$ 
        \item[(c)] Assuming that $x$ and $z$ depend on $y,$ compute
          $\dfdx{x}{y}$ and $\dfdx{z}{y}.$ 
        \item[(d)] Assume that both $x,$ $y,$ and $z$ depend on a third
          variable, $t.$ Find $\dfdx{x}{t}$ in terms of $\dfdx{y}{t},$
          and $\dfdx{z}{t}.$
        \item[(e)] Assume that both $x,$ $y,$ and $z$ depend on a third
          variable, $t.$ Find $\dfdx{y}{t}$ in terms of $\dfdx{x}{t},$
          and $\dfdx{z}{t}.$
        \item[(f)] Assume that both $x,$ $y,$ and $z$ depend on a third
          variable, $t.$ Find $\dfdx{z}{t}$ in terms of $\dfdx{x}{t},$
          and $\dfdx{y}{t}.$
        \end{description}
%         \begin{description}
%         \item[(a)] $\dxisplaystyle\dfdx{x}{y}$
%         \item[(b)] $\dxisplaystyle\dfdx{y}{x}$
%         \item[(c)] $\dxisplaystyle\dfdx{x}{z}$
%         \item[(d)] $\dxisplaystyle\dfdx{y}{x}$
%         \item[(e)] $\dxisplaystyle\dfdx{y}{z}$
%         \item[(f)] $\dxisplaystyle\dfdx{z}{x}$
%         \item[(g)] $\dxisplaystyle\dfdx{x}{z}$
%         \item[(h)] $\dxisplaystyle\dfdx{z}{y}$
%         \item[(i)] $\dxisplaystyle\dfdx{x}{x}$
%         \item[(j)] $\dxisplaystyle\dfdx{y}{y}$
%         \item[(k)] $\dxisplaystyle\dfdx{z}{z}$
%         \end{description}
  \end{embeddedproblem}

  \begin{embeddedproblem}{}
    Find $\dfdx{y}{x}$ for each of the following.
    \begin{multicols}{3}
      \begin{description}
       \item[(a)] $y=3x+2$
       \item[(b)] $y=4x-2$
       \item[(c)] $y=\frac{x}{3}+12$
       \item[(d)] $y=\frac{3x}{2}-7$
       \item[(e)] $y+3x=2$
       \item[(f)] $3x-y=0$
       \item[(g)] $ax-by=0$ ($a$ and $b$ are constants.)
       \item[(h)] $7-\frac65y+\frac85x=0$
       \item[(i)] $y=3x+2$
       \item[(j)] $1/x=5/y$
       \item[(k)] $\frac{y}{x}=\frac23$
       \item[(l)] $\frac{x}{y} -1 = \frac{2x}{3y}+2$
      \end{description}
    \end{multicols}
\end{embeddedproblem}
\begin{embeddedproblem}{}
  Two cars pass the same point at the same time. Car A is traveling at
  $30 m/h$ and Car B is traveling at $50 m/h.$ 
  \begin{description}
  \item[(a)] Use the formula
    $\text{Distance}=\text{Rate}\times\text{Time}$ to find a formula
    for Car A's position, $D_A$ after they pass.
  \item[(b)] Use the formula
    $\text{Distance}=\text{Rate}\times\text{Time}$ to find a formula
    for Car B's position, $D_B$ after they pass.
  \item[(c)] Give a formula for the distance between the cars after
    they pass.
  \item[(d)] Use the formula you found in part (c) to show that the
    cars are separating at a rate\footnote{Yes this is obvious, and
      yes, this is a lot of work just to get something so
      obvious. But, we are just beginnning and it is important that
      you learn how to use the notation of Calculus in a simple
      context before you try to use it on more complex problems. The
      problems will get harder. We promise.} of $20 m/h.$
  \item[(e)] Find the rate of change of Car A \underline{with respect
      to Car B\footnote{That is, think of Car B as fixed. What is the
        velocity of Car A?}}.
  \item[(f)] Find the rate of change of Car B \underline{with respect
      to Car A}.
  \end{description}
\end{embeddedproblem}


\end{ProblemSection}

\begin{ProblemSection}
  % \begin{embeddedproblem}{}
  %   \begin{enumerate}[label={\bf{}(\alph*)}]
  %     \begin{multicols}{2}
  %     \item $\displaystyle\dfdx{ (\inverse\sin x)}{x} = \frac{1}{\sqrt{1-x^2}} $
  %     \item $\displaystyle\dfdx{ (\inverse\cos x)}{x} =
  %       \frac{-1}{\sqrt{1-x^2}} $
  %     \item $\displaystyle\dfdx{ (\inverse\tan x)}{x} = \frac{1}{1+
  %         x^2} $
  %     \item $\displaystyle\dfdx{ (\inverse\cot x)}{x} = \frac{-1}{1+
  %         x^2}$
  %     \item $\displaystyle\dfdx{ (\inverse\sec x)}{x} =
  %       \frac{1}{\abs{x}\sqrt{x^2-1}} $
  %     \item $\displaystyle\dfdx{ (\inverse\csc x)}{x} =   \frac{-1}{\abs{x}\sqrt{x^2-1}} $      \end{multicols}
  %   \end{enumerate}
  % \end{embeddedproblem}


  
  \begin{embeddedproblem}{}
    Compute $\dfdx{y}{x}:$
    \begin{enumerate}[label={\bf{}(\alph*)}]
      \begin{multicols}{2}
      \item $\sin(x)\inverse\sin{(y)}=xy$
      \item $\inverse\cot(y) = 
        \left(
          \sqrt{x}+\frac{1}{\sqrt{x}}
        \right)\tan(y)$
      \item $\inverse\sec
        \left(
          x+y
        \right)=(x^2+y^2)$
      \item $\inverse\tan(x) =\inverse\cot(y)$
      \end{multicols}  
    \end{enumerate}
  \end{embeddedproblem}

\end{ProblemSection}


\begin{ProblemSection}
  \begin{embeddedproblem}{}
    Compute each of the following:
    \begin{enumerate}[label={\bf{}(\alph*)}]{}
    \item $\dx(e^{x^2+\sin(x)})$
    \item $\dx(e^{x}\cos(3x)$
    \item $\dx
      \left(
        \frac{3\sqrt{x}+1}{e^x}
      \right)$
    \item $\dx(\tan(e^x))$
    \end{enumerate}{}
  \end{embeddedproblem}

  \begin{embeddedproblem}{}
    For each of the    following, find an equation relating $\dx{x},$
    $\dx{y},$  $\dx{z}.$ Use this to compute $\dfdx{y}{x}$  and
    $\dfdx{y}{t}.$
    \begin{enumerate}[label={\bf{}(\alph*)}]{}
    \item $e^y=x^2+3x-2$
    \item $\sin(xy)=(e^x)^y$
    \item $\tan(x+z)=e^ye^x$
    \item $e^{x^2+y^2+z^2}=2$
    \end{enumerate}
  \end{embeddedproblem}
\end{ProblemSection}

\begin{ProblemSection}
  \begin{embeddedproblem}{}
    Solc the following equations for $x.$
    \begin{enumerate}[label={\bf{}(\alph*)}]{}
    \item $e^x=\frac{1}{\sqrt{e}}$
    \item $e^x=5$
    \item $7e^{x+3}=2$
    \item $\ln(2x-4)=0$
    \item $\ln(2x)-\ln(x+1)=4$
    \item $2\ln(x)=\ln(3x-2)$
    \end{enumerate}
  \end{embeddedproblem}
  \begin{embeddedproblem}{}
    For each of the following find $\dx{y},$ $\dfdx{y}{x},$ and
    $\dfdx{y}{t}.$
    \begin{enumerate}[label={\bf{}(\alph*)}]{}
    \item $y=\ln(2x+2)$
    \item $y=\ln
      \left(
        \frac{(2x+1)^2(3x-2)}{x^2+1}
      \right)$ \hint{Think about properties of logarithms before you
        do this.} 
    \item $\ln(y)=\sin(e^x)$
    \item $\ln(x+y)=\inverse\tan(z)$
    \end{enumerate}
  \end{embeddedproblem}
\end{ProblemSection}

\begin{ProblemSection}
  \begin{embeddedproblem}{}
    Suppose that a point $P$ is moving along the curve
    $y=x^3-x^2+x-1.$ Suppose the horizontal velocity of the point is
    constant.  Would the vertical velocity of the point be constant?
    Justify your answer.
  \end{embeddedproblem}

  \begin{embeddedproblem}{}
    Suppose that a point $P$ starts at the origin and moves along
    the positive $x$-axis with a speed of $1$ unit/sec.  At the same
    time a point $Q$ starts at the point $(1,0)$ and moves up the
    line $x=10$ with a speed of $1$ unit/sec (see below).\\
    \centerline{\includegraphics*[height=2in,width=3.5in]{../Figures/PracProb1}}
    Show that the distance between the two points, $s,$ is decreasing
    $\left(\dfdx{s}{t}<0\right)$ when $y<10-x$ and is increasing
    $\left(\dfdx{s}{t}\right)>0)$ when $y>10-x.$ Use the diagram
    above to describe what this means physically.
  \end{embeddedproblem}
\end{ProblemSection}


\begin{ProblemSection}

  \begin{embeddedproblem}{}
    \begin{enumerate}[label={\bf{}(\alph*)}]
    \item A marketing company has determined that if a concert ticket
      costs $\$50$ then $20,000$ tickets will sell.  They also determined
      that for each $\$1$ the price is increased, it will cause sales to
      go down by $500$ people.  Let $p$ be the price per ticket.  Find the
      revenue $R$ generated by the concert as a function of $p.$
    \item What should the ticket price be to maximize revenue?
    \end{enumerate}
  \end{embeddedproblem}


\end{ProblemSection}

\begin{ProblemSection}
  

  Notice that in the previous problem, the range of values for $x$ was
  $(-\infty,\infty).$  This necessitated some sort of reasoning to determine
  whether or not we have a minimum.  The hint
  suggested showing that the function was concave up.  Again, there
  are other ways to determine if the critical value provides the
  extremum desired, but there is an instance where it is relatively
  straightforward.  Consider the graph of $y=x^2$ on the interval
  $[-1,2].$
  \begin{wrapfigure}[]{r}{2in}
\captionsetup{labelformat=empty}
  \centerline{\includegraphics*[height=2in,width=2in]{../Figures/ClosedBoundedExtrema}}
\label{fig:ClosedBoundedExtrema}
\end{wrapfigure}


  Notice that $(0,0)$ is the absolute minimum of $y=x^2$ and this function
  itself has no maximum.  However, when we restrict the values of $x$ to
  the closed bounded interval $[-1,2]$ then the endpoint $(2,4)$ provides
  a maximum.  Notice that $(-1,1)$ is a local maximum as well by virtue
  of being an endpoint.  There is a theorem whose proof is beyond the
  scope of this book that says that any continuous function on a
  closed bounded interval has a maximum and a minimum.  We have not
  even rigorously defined what we mean by a continuous function yet,
  but at this point we will rely on your intuitive sense that a
  continuous function is one whose graph has no ``breaks'' in it.  A
  more precise formulation will come in the theory section.  Anyway,
  we can make use of this fact in certain optimization problems.  If
  we are trying to optimize a continuous function on a closed, bounded
  interval, then this ``Extreme Value Theorem'' guarantees that it will
  attain its maximum and its minimum somewhere on the closed, bounded
  interval.  To find it, we need only look at those values of $x$
  where the derivative is either zero or undefined, and
  the endpoints.  When we substitute these into the function,
  whichever has the largest value will be the maximum and whichever
  has the smallest value will be the minimum.  If the function is not
  continuous or if the interval is not closed or not bounded, then all
  bets are off; the function may or may not have any extrema.  In a
  physical application, we will need to determine extrema in one of
  the aforementioned ways (along with checking any endpoints).  Again,
  this emphasizes the important of determining the range of values for
  $x.$


  \begin{myexample}
    Back in Chapter\ref{chapter:ScienceBeforeCalc}, we used non-Calculus
    techniques to show that out of all rectangles with a fixed
    perimeter, the square encloses the largest area.  If you go back
    and look at the solution, it was elementary (though it did use a
    lot of high powered techniques) but it was also clever in how it
    was set up.  Let's not be so clever and solve this problem with
    our calculus ``sledgehammer.''

    Given a rectangle with sides $x$ and $y,$ the area is given by
    $A=xy$ and the perimeter is given by $P=2x+2y.$ We want to
    maximize $A$ while holding $P$ fixed.  With this in mind, we have
    $y=P/2-x$ and so
    $$
    A=A(x)=x\left(\frac{P}{2}-x\right)=\frac{P}{2} x-x^2;\ \
    0\le x\le P/2.
    $$

    Again, notice the importance of determining the range of values
    for $x.$ In this case, we have a continuous function on a closed,
    bounded interval, so it must have both a maximum and a minimum.
    We have the endpoints $x=0$ and $x=P/2.$  Setting $\dfdx{P}{x}=0,$ we have
$$ 
0=\dfdx{P}{x}=\frac{P}{2}-2x
$$
and since $\frac{P}{x}-2x$ is defined everywhere  $\left\{x=0,
  x=\frac{P}{4}, x=\frac{P}{2}\right\}$ are
the only  values of $x$ we need to consider.  Substituting into the
original formula for $A,$ we have
\begin{align*}
A(0)&=A\left(\frac{P}{2}\right)=0 \\
  A\left(\frac{P}{4}\right) &=\left(\frac{P}{4}\right)^2 \\
  &>0.
\end{align*}

Thus $x=\frac{P}{4}$ provides a maximum.  Not surprisingly, $x=0$ and
$x=\frac{P}{2}$ (so $y=0$) provide the minimum area.

To answer the question, the maximum area occurs when $x=y=\frac{P}{4}$
and so the rectangle with maximum area must be a square.
\end{myexample}

We've gotten your feet wet with some fairly straightforward
optimization problems.  It's time we apply our ``sledgehammer'' to a
more challenging problem.  Probably everyone has heard the acronym for
the colors of a rainbow {\bf{}ROY G BIV} (Red, Orange, Yellow, Green, Blue,
Indigo, Violet).  Which color is on top?  It turns out that red is on
the top of the primary rainbow.  But take a look at the pictures of
the double rainbows below.\\
\centerline{\includegraphics*[height=2in,width=5in]{../Figures/DoubleRainbow1}}
Look carefully at the  colors in the primary (inner) rainbow and
the secondary (outer) rainbow.  Do you notice that the order of the
colors is reversed?  Why do you suppose 
this happens?  There are some other questions which readily come to
mind.  Why is the primary rainbow below the secondary one?  Why is the
primary rainbow brighter and why isn't there always a secondary
rainbow?  We will use what we know about the refraction and reflection
of light to examine some of these questions.

When a ray of light enters a raindrop (which is spherical\aside{Ok,
  sure, not all raindrops are perfectly spherical. In fact, a falling
  raindrop is constantly being buffetted by the air it is passing
  through, so its shape is rather complex. Still surface tension
  forces would tend to pull it into a sphere if there were no air
  constantly bumping into it. Surface tension keeps it close enough to
sphere-shaped that, for our purposes, it doesn't matter. We will
assume that the raindrop is a sphere.}), it gets
both refracted and reflected\footnote{\textcolor{red}{We should either make sure we can use these images, or get new ones.}}.\\
\centerline{\includegraphics*[height=3in,width=5in]{../Figures/DoubleRainbow2}}
As you can see in the diagram, the ray of light entering the raindrop
at an angle of radian measure $\alpha$ is refracted to an angle of
radian measure $\beta.$  This is governed by Snell's Law
$\frac{\sin(\alpha)}{v_a} =\frac{\sin(\beta)}{v_b},$ where $v_a$ and
$v_b$ are the velocities of light in air and water, respectively.
Actually only a portion is refracted as the other portion is reflected
off the exterior of the raindrop.  Once inside the raindrop, it
reflects off of the back of the raindrop (again only a portion) so
that the angle of incidence equals the angle of reflection.  A portion
of the light then exits the raindrop governed by Snell's Law again.
The total amount of (clockwise) rotation in this process is called the
deviation angle $D(\alpha ).$

\begin{embeddedproblem}{}
 Show that $D(\alpha)=\pi-4\beta+2\alpha.$
\end{embeddedproblem}

\begin{wrapfigure}[]{r}{.5in}
\captionsetup{labelformat=empty}
\centerline{\includegraphics*[height=1.5in,width=.5in]{../Figures/DoubleRainbow3}}
\label{fig:DoubleRainbow3-ProblemSection}
\end{wrapfigure}
Notice that $\alpha$ and $\beta$ are not independent. They are
governed by Snell's Law which we will rewrite as  
$$
\sin(\alpha)=\frac{v_a}{v_b}\sin(\beta)=k \sin(\beta)
$$
where $k=\frac{v_a}{v_b}$ is the index of refraction (for water).  In
the case of white light, this index of refraction is about $4/3.$
Graphing $D(\alpha)$ for this value of $k$ and $0\le \alpha \le
\pi/2,$ we have the graph at the right.

Notice from the graph that the minimum of this curve appears to be
around the point $(1.03,2.41).$  This says that the minimum value of
$D(\alpha)$ is about $2.41 \text{ radians}\approx 138^\circ$ and it
occurs when $\alpha\approx1.03 \text{ radians}\approx 59.4^\circ.$
Notice from an enlargement the graph (below) near that minimum
that for values of $\alpha $ near $1.03$ the amount of deviation is roughly the
same.\\
\centerline{\includegraphics*[height=1in,width=5in]{../Figures/DoubleRainbow4}}
This says that light rays entering the raindrop at approximately $59^\circ$ 
have the highest concentration when exiting the raindrop.  At other
values of $\alpha$ the dispersion of light is greater.  This provides the
angle that the rainbow makes with an observer.  This supplementary
angle to $D(\alpha)$ is called the rainbow angle,
$$
R(\alpha)=\pi-D(\alpha)=\pi-(\pi-4\beta+2\alpha)=4\beta-2\alpha \text{
  radians.}
$$
For white light, we have a rainbow angle of 
$$
R(1.03)=d\pi-2.41\approx 0.732 \text{ radians}\approx 42^\circ.
$$
\centerline{\includegraphics*[height=2in,width=2in]{../Figures/DoubleRainbow5}}
To obtain the spread of colors, note that the index of refraction
changes with the color of light.  This can be seen when light shines
through a prism.\\
\centerline{\includegraphics*[height=2in,width=2in]{../Figures/DoubleRainbow6}}

\begin{embeddedproblem}{}
  \label{prob:double-rainbow-2}
  Show that for an index of refraction $k$ that $D(\alpha)$ is
  minimized when $\sin(\alpha)=\sqrt{\frac{4-k^2}{3}}.$ Use this to
  obtain $\beta=\inverse{\sin}\sqrt{\frac{4-k^2}{3k^2}}$ and fill in
  the following table.
 \begin{table}[h]
   \centering
   \begin{tabular}{|c|c|c|c|c|} \hline
     Color& k& $\alpha$ in degrees&$\beta$ in degrees&$R(\alpha)$ in
                                                       degrees\\\hline
     Red & 1.331 & &&\\\hline
     Orange & 1.332 & &&\\\hline
     Yellow & 1.333 & &&\\\hline
     Green & 1.335 & &&\\\hline
     Blue & 1.337 & &&\\\hline
     Indigo & 1.340 & &&\\\hline
     Violet & 1.344 & &&\\\hline
   \end{tabular}
 \end{table}
\end{embeddedproblem}

The values of $R(\alpha)$ you obtained in the table accounts for the
various color bands that appear in the rainbow as each angle provides
the
highest concentration of that color exiting the raindrop. \\
\centerline{\includegraphics*[height=2in,width=5in]{../Figures/DoubleRainbow7}}
The primary rainbow is created by light entering the upper portion of
the raindrop.  If light is bright enough then a secondary rainbow is
created from light entering the lower portion of the raindrop as in
the diagram below. \\
\centerline{\includegraphics*[height=2.75in,width=4in]{../Figures/DoubleRainbow8}}
This time the deviation angle $D(\alpha)$ is the total amount of
counterclockwise rotation and the rainbow angle would be
$R(\alpha)=D(\alpha)-\pi.$ Note that the extra reflection accounts for
the rainbow being not as bright.

\begin{embeddedproblem}{}
  Show that in the case of the secondary rainbow
$$
D(\alpha)=2\pi+2\alpha-6\beta
$$
and show that this is minimized when  $\sin(\alpha)=\sqrt{\frac{9-k^2}{8}}.$  Use this to show that $\beta=\inverse{\sin}\sqrt{\frac{9-k^2}{8k^2}}$ and fill in the following table.
 \begin{table}[h]
   \centering
   \begin{tabular}{|c|c|c|c|c|} \hline
     Color& k& $\alpha$ in degrees&$\beta$ in degrees&$R(\alpha)$ in
                                                       degrees\\\hline
     Red & 1.331 & &&\\\hline
     Orange & 1.332 & &&\\\hline
     Yellow & 1.333 & &&\\\hline
     Green & 1.335 & &&\\\hline
     Blue & 1.337 & &&\\\hline
     Indigo & 1.340 & &&\\\hline
     Violet & 1.344 & &&\\\hline
   \end{tabular}
 \end{table}
Use this table and the table in Problem~\ref{prob:double-rainbow-2}  to explain why the secondary
rainbow is above the primary rainbow and its colors are reversed. 
\end{embeddedproblem}
\end{ProblemSection}

\begin{ProblemSection}
  % \begin{embeddedproblem}{} %duplicate problem
  %   For each of the  IVPs below, answer these questions:
  %   \begin{enumerate}
  %   \item On what intervals is the function increasing and on what
  %     intervals is it decreasing?  Identify all local extrema.
  %   \item On what intervals is the function concave up and on what
  %     intervals is it concave down?  Identify all inflection points.
  %   \item Plot a reasonable graph for $y.$ 
  %   \end{enumerate}
  %   \begin{enumerate}[label={\bf{}(\alph*)}]
  %   \item $\dfdx{y}{t}=\sqrt{t},\ \         y(0)=1$
  %   \item $\dfdx{y}{t}=x^3-x,\ \      y(1)=3$
  %   \item $\dfdx{y}{t}=x-x^3,\ \       y(1)=3   [Compare this to the
  %     previous problem.]$
  %   \end{enumerate}
  % \end{embeddedproblem}

  \begin{embeddedproblem}{}
    Explain why a $5$th degree polynomial can have at most $4$ extrema.
    Can it have less?  Explain.
  \end{embeddedproblem}
\end{ProblemSection}

\begin{ProblemSection}
  \begin{embeddedproblem}{}
    Graph the curves, identifying all extrema and inflection points,
    given the following verbal descriptions:
    \begin{enumerate}[label={\bf{}(\alph*)}]
    \item
      \begin{multicols}{2}
        \begin{itemize}
        \item[ ] $y(0)=0$
        \item[ ]  $\left.\dfdx{y}{x}\right|_{x=-1} = 0$
        \item [ ] $\left.\dfdx{y}{x}\right|_{x= 1} = 0$
        \item [ ] $\left.\dfdxn{y}{x}{2}\right|_{x= 0} = 0$
        \item [ ] $\dfdx{y}{x} <0$ on $(-\infty,-1)\cup(1,\infty)$
        \item [ ] $\dfdx{y}{x}>0$ on $(-1,1)$
        \item [ ] $\dfdxn{y}{x}{2}>0$ on $(-\infty,0)  $
        \item [ ] $\dfdxn{y}{x}{2}<0$ on $(0,\infty)  $
        \end{itemize}
      \end{multicols}
    \item
      \begin{multicols}{2}
        \begin{itemize}
        \item [ ] $y(2)=-2$
        \item [ ] $\left.\dfdx{y}{x}\right|_{x= 2} = 1$
        \item [ ] $\left.\dfdxn{y}{x}{2}\right|_{x= 2} = 0$
        \item [ ] $\left.\dfdx{y}{x}\right|_{x= 1} \ge 0$ on
          $(-\infty, \infty)$
        \item[ ] $\dfdxn{y}{x}{2}<0$ on $(-\infty,2)$
        \item[ ] $\dfdxn{y}{x}{2}>0$ on $(2,\infty)$
        \end{itemize}
      \end{multicols}
    \item 
      \begin{multicols}{2}
        \begin{itemize}
        \item [ ] $y(0)=-2$
        \item [ ] $\left.\dfdx{y}{x}\right|_{x= 0} = 0$
        \item [ ] $\left.\dfdxn{y}{x}{2}\right|_{x=-3} = 0$
        \item [ ] $\left.\dfdxn{y}{x}{2}\right|_{x=3} = 0$
        \item [ ] $y<0$ on $(-\infty,\infty)$
        \item [ ] $\dfdx{y}{x}<0$ on $(-\infty,0)$
        \item [ ] $\dfdx{y}{x}>0$ on $(0,\infty)$
        \item [ ] $\dfdxn{y}{x}(2)<0$ on $(-\infty,-3)\cup(3,\infty)$
        \item [ ] $\dfdxn{y}{x}(2)>0$ on $(-3,3)$
        \end{itemize}
      \end{multicols}
    \item 
      \begin{multicols}{2}
        \begin{itemize}
        \item [ ] $y(0)=4$
        \item [ ] $y(-x)=y(x)$
        \item [ ] $\left.\dfdx{y}{x}\right|_{x=-2}=0$
        \item [ ] $\left.\dfdx{y}{x}\right|_{x= 0}=0$
        \item [ ] $\left.\dfdx{y}{x}\right|_{x= 2}=0$
        \item [ ] $\left.\dfdxn{y}{x}{2}\right|_{x=-2}=0$
        \item [ ] $\left.\dfdxn{y}{x}{2}\right|_{x=-1}=0$
        \item [ ] $\left.\dfdxn{y}{x}{2}\right|_{x= 1}=0$
        \item [ ] $\left.\dfdxn{y}{x}{2}\right|_{x= 2}=0$
        \item [ ] $\dfdx{y}{x}>0$ on $(-\infty,-2)\cup(-2,0)$
        \item [ ] $\dfdx{y}{x}<0$ on $(0,2)\cup(2,\infty)$
        \item [ ] $\dfdxn{y}{x}{2}<0$ on $(-\infty,-2)\cup(-1,1)\cup(2,\infty)$
        \item [ ] $\dfdxn{y}{x}{2}>0$ on $(-2,-1)\cup(1,2)$
        \end{itemize}
      \end{multicols}
    \end{enumerate}
  \end{embeddedproblem}
\end{ProblemSection}

\begin{ProblemSection}
  \begin{embeddedproblem}{}
    The following are graphs for $\dfdx{y}{x}.$  Use them to
    determine:
    \begin{enumerate}[label={\bf{}(\roman*)}]
    \item On what intervals is $y$ increasing and on what intervals is
      it decreasing?  Identify all local extrema.
    \item On what intervals is $y$ concave up and on what intervals is
      it concave down?  Identify all inflection points.
    \item Plot a reasonable graph for $y=y(x).$
  \end{enumerate}
    \begin{enumerate}[label={\bf{}(\alph*)}]
    \item\ \\ \includegraphics*[height=2in,width=2in]{../Figures/Graphing4}      
    \item\ \\ \includegraphics*[height=2in,width=2in]{../Figures/Graphing5}      
    \item\ \\ \includegraphics*[height=2in,width=2in]{../Figures/Graphing6}
  \end{enumerate}
  \end{embeddedproblem}
\end{ProblemSection}

\begin{ProblemSection}
  \begin{embeddedproblem}{}
    Graph the curves, identifying all extrema and inflection points,
    given the following verbal descriptions:
      \begin{enumerate}[label={\bf{}(\alph*)}]
      \item
        $$
        \begin{array}{ccc}
           y(0)=4&& y(-x)=y(x)\\
          \left.\dfdx{y}{x}\right|_{x=-2}=0&
          \left.\dfdx{y}{x}\right|_{x= 2}=0&
          \left.\dfdx{y}{x}\right|_{x= 0} \text{ does not exist} \\
     \left.\dfdxn{y}{x}{2}\right|_{x=-2}=0&
     \left.\dfdxn{y}{x}{2}\right|_{x=2}=0&
     \left.\dfdxn{y}{x}{2}\right|_{x=0} \text{ does not exist} \\
     \dfdx{y}{x}>0 \text{ on } (-\infty,-2)\cup(-2,0)&&
     \dfdx{y}{x}<0 \text{ on } (0,2)\cup(2,\infty)\\
     \dfdxn{y}{x}{2}<0 \text{ on }(-\infty,2)\cup(2,\infty)&&
      \dfdxn{y}{x}{2}>0 \text{ on } (-2,0)\cup(0,2)\\
        \end{array}
        $$
          % \begin{multicols}{2}
          %   \begin{list}{}{}
          %   \item [ ] 
          %   \item [ ] $$
          %   \item [ ] $
          %   \item [ ] 
          %   \item [ ] $$
          %   \item [ ] $\left.\dfdxn{y}{x}{2}\right|_{x= 1}=0$
          %   \item [ ] $\left.\dfdxn{y}{x}{2}\right|_{x= 2}=0$
          %   \item [ ] 
          %   \item [ ] $$
          %   \item [ ] 
          %   \item [ ] $$
          %   \end{list}
          % \end{multicols}
        \item           Notice that in this problem we only gave you partial
          information. You will need to figure out the rest.
          \begin{multicols}{2}
            \begin{list}{}{}
            \item[ ] $y(-x)=-y(x)$
            \item[ ] $\left.\dfdx{y}{x}\right|_{x=0}=-1$
            \item[ ] $\left.\dfdx{y}{x}\right|_{x=-3}$ does not exist
            \item[ ] $\dfdx{y}{x}>0$ on $(-\infty,-3)$
            \item[ ] $\dfdx{y}{x}<0$ on $(-3,3)$
            \end{list}
          \end{multicols}
        \item 
          \begin{multicols}{2}
            \begin{list}{}{}
            \item [ ] $y(0)=0$
            \item [ ] $\left.\dfdx{y}{x}\right|_{x=0}$ does not exist
            \item [ ] $\dfdx{y}{x}<0$ on $(-\infty,0)\cup(0,\infty)$
            \item [ ] $\dfdxn{y}{x}{2}<0$ on $(-\infty,0)$
            \item [ ] $\dfdxn{y}{x}{2}>0$ on $(0,\infty)$
            \end{list}
          \end{multicols}
        \end{enumerate}
      \end{embeddedproblem}
    \end{ProblemSection}


  \begin{embeddedproblem}{}
    Suppose that
    \begin{align*}
      \limit{x}{2}{f(x)}&=7, \\
      \limit{x}{2}{g(x)}&=3, \\
      \limit{x}{0}{\alpha(x)}&=5, \text{ and} \\
      \limit{x}{0}{\beta(x)}&=-2.
    \end{align*}
    Use the information provided above to evaluate each of the
    following limits. Cite each theorem you use as you use it. If you
    don't have enough information, state precisely what else is
    needed.
    \begin{multicols}{2}
      \begin{enumerate}[label={\bf{}(\alph*)}]
      \item $\limit{x}{2}{(f(x)+g(x)}$
      \item $\limit{x}{2}{4f(x)}$
      \item $\limit{x}{2}{-2g(x)}$
      \item $\limit{x}{2}{\left(-5f(x)+\frac13g(x)\right)}$
      \item $\limit{x}{0}{\left(-5f(x)+\frac13g(x)\right)}$
      \item $\limit{t}{2}{\left(-5f(t)+\frac13g(t)\right)}$
      \item $\limit{s}{0}{\alpha(s)}$
      \item $\limit{x}{2}{\left(f(x)+\alpha(x)\right)}$
      \item $\limit{s}{0}{\left(\alpha(s)+\beta(s)\right)}$
      \item $\limit{t}{2}{\left(f(t)+g(t)+\alpha(t-2)+\beta(t-2)\right)}$
      \end{enumerate}
    \end{multicols}
  \end{embeddedproblem}
  \begin{embeddedproblem}{}
    \begin{multicols}{2}{Compute each of the following limits:}
      \begin{enumerate}[label={\bf{}(\alph*)}]
        \item $\limit{x}{3}{\frac{x^2-9}{x-3}}$
        \item $\limit{x}{4}{\frac{x^2-16}{x-4}}$
        \item $\limit{x}{\sqrt{2}}{\frac{x^2-2}{x-\sqrt{2}}}$
        \item $\limit{x}{a}{\frac{x^2-a^2}{x-a}}$
        \end{enumerate}
      \end{multicols}
    \end{embeddedproblem}
    
    \begin{embeddedproblem}{}
      Compute each of the following limits:
      \begin{multicols}{2}
        \begin{enumerate}[label={\bf{}(\alph*)}]
        \item $\limit{x}{1}{\frac{x^3-1}{x-1}}$
        \item $\limit{x}{2}{\frac{x^3-8}{x-2}}$
        \item $\limit{x}{3}{\frac{x^3-27}{x-3}}$
        \item $\limit{x}{5}{\frac{x^3-125}{x-5}}$
        \item $\limit{x}{a}{\frac{x^3-a^3}{x-a}}$
        \item $\limit{x}{1}{\frac{x^3+1}{x+1}}$
        \item $\limit{x}{2}{\frac{x^3+8}{x+2}}$
        \item $\limit{x}{3}{\frac{x^3+27}{x+3}}$
        \item $\limit{x}{5}{\frac{x^3+125}{x+5}}$
        \item $\limit{x}{a}{\frac{x^3+a^3}{x+a}}$
        \end{enumerate}
      \end{multicols}
    \end{embeddedproblem}
    
    \begin{embeddedproblem}{}
      Compute each of the following limits:
      \begin{multicols}{2}
        \begin{enumerate}[label={\bf{}(\alph*)}]
        \item $\limit{x}{a}{\frac{x^4-a^4}{x-a}}$
        \item $\limit{x}{a}{\frac{x^5-a^5}{x-a}}$
        \item $\limit{x}{a}{\frac{x^n-a^n}{x-a}}$
        \item $\limit{x}{a}{\frac{x^4+a^4}{x+a}}$
        \item $\limit{x}{a}{\frac{x^5+a^5}{x+a}}$
        \item $\limit{x}{a}{\frac{x^n+a^n}{x+a}}$
        \end{enumerate}
      \end{multicols}
    \end{embeddedproblem}
    
    \begin{embeddedproblem}{}
      Compute each of the following limits:
      \begin{multicols}{2}
        \begin{enumerate}[label={\bf{}(\alph*)}]
        \item $\limit{x}{2}{\frac{x^2+x-6}{x-2}}$
        \item $\limit{x}{-3}{\frac{x^2+x-6}{x+3}}$
        \item $\limit{x}{1}{\frac{x^5-x^4+x^3-x^2}{x-1}}$
        \end{enumerate}
      \end{multicols}
    \end{embeddedproblem}
    
    It should be clear that  in all of the problems above the goal
    was to manipulate the given expression algebraically until we can
    divide out the troublesome factor in the denominator.

    if we're asked to evaluate, say
    $\limit{x}{1}{\frac{\inverse{x}-1}{x-1}},$ that goal remains but
    the algebra changes:
    \begin{align*}
      \limit{x}{1}{\frac{\inverse{x}-1}{x-1}} &=
                                                \limit{x}{1}{\frac{\frac1x-1}{x-1}} \\
                                              &=
                                                \limit{x}{1}{\frac{\frac{x-1}{x}}{x-1}} \\
                                              &=\limit{x}{1}{\frac{x-1}
                                                {x(x-1)}} \\
      &=\limit{x}{1}{\frac1x} \\
      &=1.
    \end{align*}
    \begin{embeddedproblem}{}
      Compute each of the following limits:
      \begin{multicols}{2}
        \begin{enumerate}[label={\bf{}(\alph*)}]
          \item $\limit{x}{1/2}{\frac{\inverse{x}-2}{x-1/2}}$
          \item $\limit{x}{1/3}{\frac{\inverse{x}-3}{x-1/3}}$
          \item $\limit{x}{2}{\frac{\inverse{x}-1/2}{x-2}}$
          \item $\limit{x}{a}{\frac{\inverse{x}-a}{x-1/a}}$
          \end{enumerate}
        \end{multicols}
      \end{embeddedproblem}

      Another variation on this theme is roots (fractional exponents).

      \begin{embeddedproblem}{}
      Compute each of the following limits:
      \begin{multicols}{2}
        \begin{enumerate}[label={\bf{}(\alph*)}]
          \item $\limit{x}{4}{\frac{x^{1/2}-2}{x-4}}$
          \item $\limit{x}{9}{\frac{x^{1/2}-3}{x-9}}$
          \item $\limit{x}{16}{\frac{x^{1/2}-4}{x-16}}$
          \item $\limit{x}{a^2}{\frac{x^{1/2}-a}{x-a^2}}$
          \end{enumerate}          
        \end{multicols}
      \end{embeddedproblem}

    
    
    
\begin{ProblemSection}
  \begin{embeddedproblem}{}
    Determine the following limits.  Don't forget that your answer
    should be a statement not just a number.
    \begin{multicols}{2}
      \begin{enumerate}[label={\bf{}(\alph*)}]
      \item $\limit{x}{-\infty}{\frac{x^2+1}{x^3}}$
      \item $\limit{x}{ \infty}{\frac{x-2}{x^2-4}}$
      \item $\limit{x}{ \infty}{\cos\left(\frac1x\right)}$
      \item $\limit{x}{-\infty}{\cos\left(\frac1x\right)}$
      \item $\limit{x}{ \infty}{\cos\left(\inverse\tan(x)\right)}$
      \item $\limit{x}{-\infty}{\cos\left(\inverse\tan(x)\right)}$
      \end{enumerate}
    \end{multicols}
  \end{embeddedproblem}
    The limit of a quotient works the same as product, with the
    proviso that you cannot divide by zero.  Specifically, if
    $\limit{x}{\infty}{f(x)}=L$ and $\limit{x}{\infty}{g(x)}=M\ne0$
    then
    $$
    \limit{x}{\infty}{\frac{f(x)}{g(x)}}=\frac{L}{M}.
    $$
    \begin{embeddedproblem}{}
      Determine the following limits.  Don't forget that your answer
      should be a statement not just a number.
      \begin{multicols}{2}
        \begin{enumerate}[label={\bf{}(\alph*)}]
        \item $\limit{x}{\infty}{\frac{\inverse\tan(x)}{1+\frac1x}}$
        \item $\limit{x}{-\infty}{\frac{\cos\left(\frac1x\right)}{x}}$
        \end{enumerate}
      \end{multicols}
    \end{embeddedproblem}
\end{ProblemSection}

  \begin{ProblemSection}
    \begin{embeddedproblem}{}
      \begin{multicols}{2}[    Determine each of the following
        limits. If a limit does not exist, then say so.]
        \begin{enumerate}[label={\bf{}(\alph*)}]
        \item $\limit{x}{\infty}{\frac{3x^4-20x^2+100}{4x^4-x}}$
        \item $\limit{x}{-\infty}{\frac{3x^5+20x^4+100}{4x^4+x^2}}$
        \item $\limit{x}{\infty}{\frac{2x^{\frac32}+x^{\frac12}+100}{x^{\frac32}+x}}$
        \item $\limit{x}{\infty}{\frac{2x^{\frac32}+x^{\frac12}+100}{x^{\frac32}+x^2}}$
        \item $\limit{x}{\infty}{\sqrt{x^2+5x+1}-x}$
        \item $\limit{x}{-\infty}{\sqrt{x^2+5x+1}-x}$
        \end{enumerate}
      \end{multicols}
  \end{embeddedproblem}
  \end{ProblemSection}

\begin{ProblemSection}
  \begin{embeddedproblem}{}
    \begin{enumerate}[label={\bf{}(\alph*)}]
    \item $\limit{x}{\infty}{\frac{x\cos(x)}{x^2+1}}$
    \item $\limit{x}{-\infty}{\frac{x\cos(x)}{x^2+1}}$ (Be careful
      with the inequalities!)
    \item $\limit{x}{\infty}{\frac{\floor{x}}{x^2}}$ where $\floor{x}$
      represents the greatest integer less than or equal
      to\footnote{This is sometimes called the ``round down''
        function, or the ``step'' function.  For example,
        $\floor{5.7}=5,$ $\floor{2}=2,$ $\floor{\pi}=3,$
        $\floor{-5.7}=-6,$ etc.} $x.$
    \item 	Suppose $\abs{f(x)}\le B$ for some fixed
      number\footnote{We say that f(x) is bounded.} $B.$   Show that
      if $\limit{x}{\infty}{g(x)}=0,$ then
      $\limit{x}{\infty}{f(x)\cdot g(x)}=0.$  [Note: You cannot just
      say $\limit{x}{\infty}{f(x)\cdot g(x)}
      = \limit{x}{\infty}{f(x)}\cdot\limit{x}{\infty}{g(x)}.$ Why not?
    \end{enumerate}
  \end{embeddedproblem}
\end{ProblemSection}

\begin{ProblemSection}
  \begin{embeddedproblem}{}
    Determine the following limits. If a limit does not exist, then
    say so.
    \begin{enumerate}
    \item
      \begin{enumerate}
      \item[$(\ast)$]
$\llimit{x}{1}{\frac{1}{1-x}}$
        % $\displaystyle\underset{x<1}{\lim_{x\rightarrow1}}\frac{1}{1-x}$
\item[$(\ast\ast)$]
%  $\rlimit{x}{1}{\frac{1}{1-x}}$
        $\displaystyle\underset{x>1}{\lim_{x\rightarrow1}}\frac{1}{1-x}$
      \item[$(\ast\ast\ast)$]
        $\limit{x}{1}{\frac{1}{1-x}}$
      \end{enumerate}
    \item
      \begin{enumerate}
      \item[$(\ast)$]
        $\llimit{x}{0}{\csc{x}}$
      \item[$(\ast\ast)$]
        $\rlimit{x}{0}{\csc{x}}$
      \item[$(\ast\ast\ast)$]
        $\limit{x}{0}{\csc(x)}$
      \end{enumerate}
    \item
      \begin{enumerate}
      \item[$(\ast)$]
        $\llimit{x}{1}{\frac{1}{(1-x)^{\frac23}}}$
      \item[$(\ast\ast)$]
        $\rlimit{x}{1}{\frac{1}{(1-x)^{\frac23}}}$
      \item[$(\ast\ast\ast)$]
        $\limit{x}{1}{\frac{1}{(1-x)^{\frac23}}}$
      \end{enumerate}
    \end{enumerate}
  \end{embeddedproblem}

  \begin{embeddedproblem}{}
    	Determine all horizontal and vertical asymptotes for the
        following curves. For the horizontal asymptotes, make sure to
        indicate whether the limits are as $x$ approaches infinity or
        negative infinity.  For the vertical asymptotes, make sure to
        indicate whether they are one sided or two sided limits.
        \begin{multicols}{3}
          \begin{enumerate}[label={\bf{}(\alph*)}]
          \item $y=\dfrac{x}{x-2}$
          \item $y=\dfrac{x}{(x-2)^2}$
          \item $y=\dfrac{x-2}{x^2-4}$
          \end{enumerate}
        \end{multicols}
      \end{embeddedproblem}
\end{ProblemSection}

\begin{ProblemSection}
  \begin{embeddedproblem}{}
    \begin{multicols}{2}[Determine each of the following limits.]
      \begin{enumerate}[label={\bf{}(\alph*)}]
      \item $\limit{x}{3}{\frac{x^3-27}{x-3}}$
      \item $\limit{x}{3}{\frac{x^3-27}{x^2-9}}$
      \item $\limit{h}{0}{\frac{(2+h)^2-4}{h}}$
      \item $\limit{h}{0}{\frac{\frac{1}{1+h}-1}{h}}$
      \item $\limit{x}{4}{\frac{\sqrt{x}-2}{x-4}}$
      \end{enumerate}
    \end{multicols}
  \end{embeddedproblem}
\end{ProblemSection}

\begin{ProblemSection}
  \begin{embeddedproblem}{}
    \begin{multicols}{3}[Determine each of the following limits. If a
      limit does not exist, say so.]
      \begin{enumerate}[label={\bf{}(\alph*)}]
      \item $\limit{x}{-3}{\frac{x^2-9}{x^3-27}}$
      \item $\limit{x}{0}{\frac{x}{\ln(x+1)}}$
      \item $\displaystyle\underset{x>1}{\lim_{x\rightarrow 1}}\frac{\sqrt[3]{x-1}}{\sqrt{x-1}}$
      \item $\limit{x}{1}{\frac{\sqrt[3]{x-1}}{\sqrt{x-1}}}$
      \item $\limit{x}{0}{\frac{e^{5x}-1}{x-1}}$
      \item $\limit{x}{0}{
          \left(
            \frac1x-\frac{1}{e^x-1}
          \right)}$
      \item $\limit{x}{0}{\frac{e^x-1}{\tan(2x)}}$
      \item $\limit{x}{1}{\frac{1-\sqrt{x}}{\ln(x)}}$
      \item $\displaystyle\underset{x>1}{\lim_{x\rightarrow 1}}\frac{\ln(x)}{\sqrt{x}}$
      \end{enumerate}
    \end{multicols}
  \end{embeddedproblem}
\end{ProblemSection}

  \begin{ProblemSection}
    \begin{embeddedproblem}{}
      \begin{multicols}{2}
        \begin{enumerate}[label={\bf{}(\alph*)}]
        \item $\displaystyle\underset{x>0}{\lim_{x\rightarrow 0}}\frac{\ln(x)}{\inverse{x}}$
        \item $\displaystyle\underset{x>0}{\lim_{x\rightarrow 0}}x\ln(x)$
        \item $\limit{x}{\infty}{\frac{\ln(x)}{\csc{x}}}$
        \item $\displaystyle\underset{x>0}{\lim_{x\rightarrow 0}}\left(\pi/2-x\right)\cdot\tan(x)$
        \item $\limit{x}{-\infty}{x\cdot e^x}$
        \end{enumerate}
      \end{multicols}
    \end{embeddedproblem}
  \end{ProblemSection}

    \begin{ProblemSection}
      \begin{embeddedproblem}{}
        \begin{multicols}{3}[Determine each of the following limits. If a limit doesn't exist, then explain why not.]
          \begin{enumerate}[label={\bf{}(\alph*)}]
          \item $\limit{x}{2}{}{\frac{\frac1x-\frac12}{x-2}}$
          \item $\limit{x}{0}{\frac{e^x-e^{-x}}{x}}$
          \item $\limit{x}{\infty}{\frac{4x+\ln(x)}{x+6}}$
          \item $\limit{x}{-\infty}{x^2e^x}$
          \item $\limit{x}{\infty}{\frac{x^3+2}{x^3-1}}$
          \item $\limit{x}{0}{(e^x-3x)^{\frac1x}}$
          \item $\llimit{x}{<2}{2}{\frac{x^2+4}{x^2-4}}$
          \item $\limit{x}{1}{
              \left(
                \frac{1}{\ln(x)}-\frac{x}{x-1}
              \right)}$
          \item $\llimit{x}{<\frac{\pi}{2}} {\frac{\pi}{2}} {(\tan(x))^{\cos(x)}}$
          \item
            $\llimit{x}{>\frac{\pi}{2}} {\frac{\pi}{2}} {(\tan(x))^{\cos(x)}}$
          \item $\limit{x}{\frac{\pi}{2}}{(\tan(x))^{\cos(x)}}$
          \item $\limit{x}{\infty}{
              \left(
                \frac{\ln(x)}{x}-\frac{1}{\sqrt{x}}
              \right)}$
          \end{enumerate}
        \end{multicols}
      \end{embeddedproblem}
      \begin{embeddedproblem}{}
        [Taken from Shaum's Outline on Calculus.]  The current $i$ in
        a coil with re qsistance $R,$ inductance $L,$ and a constant
        electromotive force $E,$ at time $t$ is given by
        $$
        i=\frac{E}{R}\left(1-e^{-\frac{Rt}{L}} \right).
        $$
        Obtain an estimate for $i$ as the resistance approaches
        zero. (This is called superconductivity and typically involves
        supercooling)
      \end{embeddedproblem}
      \begin{embeddedproblem}{}
        [Essentially taken from Stewart.]  In a paper on the
        sensitivity of the retina published in $1933,$ W.~Stanley
        Stiles and B.~H.~Crawford determined that the percentage $P$
        of total luminance entering an eye's pupil which is sensed at
        the retina can be modeled by
        $$
        P=\frac{1-10^{-\rho{}r^2}}{\rho{}r^2  \ln(10)}
        $$
        where $r$ is the radius of the pupil, and $\rho{}$ is an experimentally
        determined constant, typically about $0.05.$
        \begin{enumerate}[label={\bf{}(\alph*)}]
        \item What is the percentage of luminance sensed by a pupil of
          radius $3$mm, assuming that $\rho=0.05?$
        \item What is the percentage of luminance sensed by a pupil of
          radius $2$mm, assuming that
          $\rho=0.05?$ Is this value larger than that in part (a)?
          Does this make sense?
        \item Determine
          $\llimit{r}{r>0}{0}{P}$ Does this answer make sense?  Would it
          ever be attained?
        \end{enumerate}
      \end{embeddedproblem}
      \begin{embeddedproblem}{}
        Recall that we talked about the growth rate of a population of
        fish (measured in tons) in a lake and modeled it with the
        differential equation
        $$
        \dfdx{y}{t}=.01y\left(1-\frac{y}{100}\right)
        $$
        What is the physical difference between saying
        $\left.\dfdx{y}{t}\right|_{y=100}=0$ and
        $\llimit{y}{<100}{100}{\dfdx{y}{t}}=0?$
      \end{embeddedproblem}
    \end{ProblemSection}

\begin{ProblemSection}
  \begin{embeddedproblem}{}
    Find all values of $x$ for which each of the following curves is
    increasing and for which it is decreasing. \\
         \hint{
         Recall that the definition of the absolute value is: 
         $$
         \abs{x} =
         \begin{cases}
           x& \text{if } x\ge0 \\
           -x& \text{if } x<0 
         \end{cases}.
         $$}
    \begin{multicols}{3}
      \begin{description}
       \item[(a)] $y=\frac{1}{x^2}$
       \item[(b)] $y=\frac{x^2}{x+1}$
       \item[(c)] $y=\abs{x}$
       \item[(d)] $y=\cos^2x -\cos x, 0\le x\le 1\pi$
       \item[(e)] $y=$
       \item[(f)] $y=$
       \item[(g)] $y=$
       \item[(h)] $y=$
       \item[(i)] $y=$
       \item[(j)] $y=$
       \item[(k)] $y=$
       \item[(l)] $y=$
      \end{description}
    \end{multicols}
\end{embeddedproblem}

  \begin{embeddedproblem}{} 
    Consider the top view of an airplane $P$ starting at the point
    $(1,0)$ and traveling up the line $x=1$ at a constant speed $v.$
    When the plane is at $(1,0),$ a homing missile $M$ is fired from
    the origin directly at the plane. Assuming that the missile
    travels at a speed which is $k$ times the speed of the plane
    $(k>1)$ and is always aimed directly at the plane, find the path
    the missile takes.  [Such a path is called a pursuit curve.]  The
    diagram below shows the situation at time $t.$\\
    \centerline{\includegraphics*[height=3in,width=4in]{../Figures/PursuitCurve}}
    \begin{description}
    \item [(a)]  If we let $s$ denote the distance the missile has traveled
      at time $t,$ show that the pursuit curve must satisfy the
      differential equation
      \[\dfdx{y}{x}=\frac{(\frac{s}{k}-y)}{1-x}\]
      and use this to show that the pursuit curve must satisfy the
      differential equation
      \[(1-x) \dfdxn{y}{x}{2}=1/k \dfdx{s}{x}=1/k
        \sqrt{1+(\dx{y}/\dx{x})^2}\]
      with the initial conditions $y(0)=0,$ $\dfdx{y}{x}(0)=0.$
    \item[(b)]  Solve the differential equation in part (a).  \hint{Let
      $z=\dfdx{y}{x}$ and solve the resulting equations for $z.$}
    \item[(c)]  How long does it take for the missile to catch the plane?
    \end{description}
  \end{embeddedproblem}

\begin{embeddedproblem}{}
  \begin{description}
  \item[(a)] Show that $x^2=1+2(x-1) +(x-1)^2.$
  \item[(b)] Find the equation of the straight line tangent to the
    graph of $y=x^2$ at the point $(1,1).$
  \item[(c)] How is your formula in part (b) related to the formula in
    part (a)?
  \end{description}
\end{embeddedproblem}

\begin{embeddedproblem}{}
  Consider the functions:\marginpar{This belongs with the foreshadowing of power series}
\[    y(x) = \frac{x^2}{x^2+1} \text{ and } p(x)=Ax^2+Bx+C.\] 
We want to find values for the parameters, $A, B,$ and $C$ that make
the graphs of these two functions as similar as possible near the
point $(2, 4/5).$
  \begin{description}
  \item[(a)] First compute $\dfdx{y}{x},$ and $\dfdxn{y}{x}{2}$ at the
    point $x=2.$
  \item[(b)] Naturally we'll want both functions to pass through the
    point $(2, 4/5)$ so we set 
    \begin{align}
      p(2) &= f(2).\label{TS:1}
      \intertext{To guarantee that the graphs are ``pointing in the
        same direction'' we set their derivatives at the point $(2,
        4/5)$ equal as well.}
      y^\prime(2) &= p^\prime(2)\label{TS:2}
      \intertext{To guarantee that the graphs are ``similarly curved''
        we set their second derivatives at the point $(2,
        4/5)$ equal as well.}
      y^{\prime\prime}(2) &= p^{\prime\prime}(2)\label{TS:3}
    \end{align}
    Solve equations~\ref{TS:1},~\ref{TS:2}, and~\ref{TS:3} to find $A,
    B,$ and $C.$
  \item[(c)] On the same set of axes graph $f(x)$ and $p(x).$ How are
    these graphs similar and how are they different?
  \end{description}

\end{embeddedproblem}

    \begin{embeddedproblem}{}
      \marginpar{This should be somewhere else.}  L'Hopital's Rule
      problem: $\limit{x}{0}{(1-ax)^{\frac1x}}.$ (Form: $0^0$ but is
      \underline{not} equal to $1.$
  \end{embeddedproblem}

\end{ProblemSection}

\begin{ProblemSection}
  \begin{embeddedproblem}{}
    The area of a rectangle is $400 m^2.$  Find a formula for the
    perimeter of the rectangle in terms of one of the sides of the
    rectangle.  Remember to define what your variables are (along with
    their units).  Also remember to include the possible values for
    your independent variable (the domain of your function).
  \end{embeddedproblem}
  \begin{embeddedproblem}{}
    \begin{description}
    \item[a)] You throw a ball straight up with an initial velocity of
      $50 m/s.$ Assuming that there is no air resistance and that the
      ball accelerates toward the earth at a constant rate of
      $9.8 ((m/s))/s,$ what would the ball's velocity be at $t$
      seconds.
   \item[b)] How would your formula change if you threw the ball
     downward?
    \end{description}
  \end{embeddedproblem}
  \begin{embeddedproblem}{}
    If a promoter charges $\$50$ each for concert tickets, then he/she
    will sell $15000$ tickets.  For each $\$1$ increase in price per
    ticket, $100$ tickets less will be sold.  Determine a formula for
    the revenue based on the price per ticket, assuming that the price
    will be at least $\$50$ per ticket.
  \end{embeddedproblem}
  \begin{embeddedproblem}{}
    The perimeter of an isosceles triangle is $100$ units.  Find the
    area of the triangle as a function of the base.  Don't forget to
    define your variables and give the domain of the function.
  \end{embeddedproblem}
  \begin{embeddedproblem}{}
    A cylindrical tank must hold $500$ cubic feet of water.  The ends of
    the tank are cut from two individual squares of metal, and the
    side is constructed from one rectangular sheet of metal.  The
    metal for the ends costs $\$5$ per square foot and the metal for the
    side costs $\$4$ per square foot.  The weld for the seams costs $\$2$
    per linear foot.  What is the cost of the tank (a) As a function of
    its height?  (b)  As a function of the radius?  Don't forget
    to identify variables and to provide the domain.
  \end{embeddedproblem}
  \begin{embeddedproblem}{}
    Consider a rectangular box with a square base whose volume is
    $500~ m^3.$ Find the surface area of the box as a function of its
    height.  Now find it as a function of the length of the side of
    the base.  Don't forget to identify your variables and to provide
    the domain.
  \end{embeddedproblem}
  \begin{embeddedproblem}{}
    \begin{description}
    \item[(a)] The fixed cost for publishing a specific Calculus book
      is $\$100,000.$ These are costs that the publisher would incur
      no matter how many books were printed.  Beyond that, it costs
      $\$20$ per book to actually print the book.  If we let n denote
      the number of books printed and $C=C(n)$ denote the cost in
      dollars from printing this many books, then give a formula for
      the total cost to print $n$ books.  What would the cost per book
      be?
    \item[(b)] Suppose the company plans to charge $\$150$ per book.
      Assuming the company can sell as many books as it produces (as
      with a print-on-demand company), write a formula for the profit
      $P$ in dollars that the company would have from printing and
      selling $n$ books.  How many books would need to be sold to
      break even?
    \end{description}
  \end{embeddedproblem}
\end{ProblemSection}

\begin{ProblemSection}
  \begin{embeddedproblem}{}
    Suppose that
    \begin{align*}
      \limit{x}{2}{f(x)}&=7, \\
      \limit{x}{2}{g(x)}&=3, \\
      \limit{x}{0}{\alpha(x)}&=5, \text{ and} \\
      \limit{x}{0}{\beta(x)}&=-2.
    \end{align*}
    Use the information provided above to evaluate each of the
    following limits. Cite each theorem you use as you use it. If you
    don't have enough information, state precisely what else is
    needed.
    \begin{multicols}{2}
      \begin{enumerate}[label={\bf{}(\alph*)}]
      \item $\limit{x}{2}{(f(x)+g(x)}$
      \item $\limit{x}{2}{4f(x)}$
      \item $\limit{x}{2}{-2g(x)}$
      \item $\limit{x}{2}{\left(-5f(x)+\frac13g(x)\right)}$
      \item $\limit{x}{0}{\left(-5f(x)+\frac13g(x)\right)}$
      \item $\limit{t}{2}{\left(-5f(t)+\frac13g(t)\right)}$
      \item $\limit{s}{0}{\alpha(s)}$
      \item $\limit{x}{2}{\left(f(x)+\alpha(x)\right)}$
      \item $\limit{s}{0}{\left(\alpha(s)+\beta(s)\right)}$
      \item $\limit{t}{2}{\left(f(t)+g(t)+\alpha(t-2)+\beta(t-2)\right)}$
      \end{enumerate}
    \end{multicols}
  \end{embeddedproblem}
  \begin{embeddedproblem}{}
    \begin{multicols}{2}{Compute each of the following limits:}
      \begin{enumerate}[label={\bf{}(\alph*)}]
        \item $\limit{x}{3}{\frac{x^2-9}{x-3}}$
        \item $\limit{x}{4}{\frac{x^2-16}{x-4}}$
        \item $\limit{x}{\sqrt{2}}{\frac{x^2-2}{x-\sqrt{2}}}$
        \item $\limit{x}{a}{\frac{x^2-a^2}{x-a}}$
        \end{enumerate}
      \end{multicols}
    \end{embeddedproblem}
    
    \begin{embeddedproblem}{}
      Compute each of the following limits:
      \begin{multicols}{2}
        \begin{enumerate}[label={\bf{}(\alph*)}]
        \item $\limit{x}{1}{\frac{x^3-1}{x-1}}$
        \item $\limit{x}{2}{\frac{x^3-8}{x-2}}$
        \item $\limit{x}{3}{\frac{x^3-27}{x-3}}$
        \item $\limit{x}{5}{\frac{x^3-125}{x-5}}$
        \item $\limit{x}{a}{\frac{x^3-a^3}{x-a}}$
        \item $\limit{x}{1}{\frac{x^3+1}{x+1}}$
        \item $\limit{x}{2}{\frac{x^3+8}{x+2}}$
        \item $\limit{x}{3}{\frac{x^3+27}{x+3}}$
        \item $\limit{x}{5}{\frac{x^3+125}{x+5}}$
        \item $\limit{x}{a}{\frac{x^3+a^3}{x+a}}$
        \end{enumerate}
      \end{multicols}
    \end{embeddedproblem}
    
    \begin{embeddedproblem}{}
      Compute each of the following limits:
      \begin{multicols}{2}
        \begin{enumerate}[label={\bf{}(\alph*)}]
        \item $\limit{x}{a}{\frac{x^4-a^4}{x-a}}$
        \item $\limit{x}{a}{\frac{x^5-a^5}{x-a}}$
        \item $\limit{x}{a}{\frac{x^n-a^n}{x-a}}$
        \item $\limit{x}{a}{\frac{x^4+a^4}{x+a}}$
        \item $\limit{x}{a}{\frac{x^5+a^5}{x+a}}$
        \item $\limit{x}{a}{\frac{x^n+a^n}{x+a}}$
        \end{enumerate}
      \end{multicols}
    \end{embeddedproblem}
    
    \begin{embeddedproblem}{}
      Compute each of the following limits:
      \begin{multicols}{2}
        \begin{enumerate}[label={\bf{}(\alph*)}]
        \item $\limit{x}{2}{\frac{x^2+x-6}{x-2}}$
        \item $\limit{x}{-3}{\frac{x^2+x-6}{x+3}}$
        \item $\limit{x}{1}{\frac{x^5-x^4+x^3-x^2}{x-1}}$
        \end{enumerate}
      \end{multicols}
    \end{embeddedproblem}
    
    It should be clear that  in all of the problems above the goal
    was to manipulate the given expression algebraically until we can
    divide out the troublesome factor in the denominator.

    if we're asked to evaluate, say
    $\limit{x}{1}{\frac{\inverse{x}-1}{x-1}},$ that goal remains but
    the algebra changes:
    \begin{align*}
      \limit{x}{1}{\frac{\inverse{x}-1}{x-1}} &=
                                                \limit{x}{1}{\frac{\frac1x-1}{x-1}} \\
                                              &=
                                                \limit{x}{1}{\frac{\frac{x-1}{x}}{x-1}} \\
                                              &=\limit{x}{1}{\frac{x-1}
                                                {x(x-1)}} \\
      &=\limit{x}{1}{\frac1x} \\
      &=1.
    \end{align*}
    \begin{embeddedproblem}{}
      Compute each of the following limits:
      \begin{multicols}{2}
        \begin{enumerate}[label={\bf{}(\alph*)}]
          \item $\limit{x}{1/2}{\frac{\inverse{x}-2}{x-1/2}}$
          \item $\limit{x}{1/3}{\frac{\inverse{x}-3}{x-1/3}}$
          \item $\limit{x}{2}{\frac{\inverse{x}-1/2}{x-2}}$
          \item $\limit{x}{a}{\frac{\inverse{x}-a}{x-1/a}}$
          \end{enumerate}
        \end{multicols}
      \end{embeddedproblem}

      Another variation on this theme is roots (fractional exponents).

      \begin{embeddedproblem}{}
      Compute each of the following limits:
      \begin{multicols}{2}
        \begin{enumerate}[label={\bf{}(\alph*)}]
          \item $\limit{x}{4}{\frac{x^{1/2}-2}{x-4}}$
          \item $\limit{x}{9}{\frac{x^{1/2}-3}{x-9}}$
          \item $\limit{x}{16}{\frac{x^{1/2}-4}{x-16}}$
          \item $\limit{x}{a^2}{\frac{x^{1/2}-a}{x-a^2}}$
          \end{enumerate}          
        \end{multicols}
      \end{embeddedproblem}
    \end{ProblemSection}

\begin{ProblemSection}
  \begin{embeddedproblem}{}
    In this section we have shown that\footnote{We've actually beaten
      this problem nearly to death.}
    $$\limit{x}{\infty}{\frac{1}{x}}=0$$
    Use that fact to show that each of the following is also equal to
    zero.
    \begin{multicols}{2}
      \begin{enumerate}[label={\bf{}(\alph*)}]
        \item $\limit{x}{\infty}{}$
        \item $\limit{x}{\infty}{}$
        \item $\limit{x}{\infty}{}$
        \item $\limit{x}{\infty}{}$
        \item $\limit{x}{\infty}{}$
      \item 
      \end{enumerate}
    \end{multicols}
  \end{embeddedproblem}

  \begin{embeddedproblem}{}
      Prove that  $\limit{x}{\infty}{1+\frac{1}{x+1}}=1$ and use that
      to show that each of the following limits is also equal to one.
      \begin{multicols}{2}
        \begin{enumerate}[label={\bf{}(\alph*)}]
        \item $\limit{x}{\infty}{2-\left(\frac{x}{x+1}\right)}$
        \item $\limit{x}{\infty}{\frac{x+2}{x+1}}$
        \item $\limit{x}{\infty}{\frac{x^2+3x+2}{(x+1)^2}}$
        \end{enumerate}
      \end{multicols}
    \end{embeddedproblem}
  \end{ProblemSection}

  \begin{embeddedproblem}{}
     \textcolor{red}{This should become a problem in the "theory"
    section.}


    The classic example of this is the absolute value function:
    $y=\abs{x}$ whose graph looks like this:
    \centerline{\includegraphics*[height=2in,width=2in]{../Figures/AbsVal}}
    At the point $(0,0)$ we cannot identify a unique tangent line, so
    we say that the derivative of the graph of $y=\abs{x}$ \emph{does
      not exist} at the point $(0,0).$ The reasons for this are
    actually pretty deep and get at the heart of what makes Calculus
    work. For that reason we'd really prefer to not even consider
    functions without derivatives right now. Unfortunately, they do
    come up rather naturally, so we must find a way to deal with them
    even if we don't fully understand them yet.
  \end{embeddedproblem}
%%%%%%%%%%END

%%% Local Variables: 
%%% mode: latex
%%% outline-minor-mode: t
%%% eval:(flyspell-mode)
%%% TeX-master: "DiffCalc"
%%% End: 
